% =============================================================================
%  CHAPTER: BACKGROUND
% =============================================================================
%  Master's thesis - Extension of "Balancing KPIs through Multi-Objective
%  Prescriptive Process Analytics" (Le, Buliga, Ronzani, Padella, de Leoni).
%
%  Thesis contribution: quantification of the predictive uncertainty
%  (aleatoric + epistemic) of the two CatBoost models, and use of "confidence
%  as a KPI" in the Pareto front, together with probability calibration and
%  diversified top-k selection.
%
%  HOW TO USE THIS FILE
%  -------------------
%  - This is only the skeleton of the chapter: every section/subsection is a
%    real LaTeX command; ALL the guidance text is in comments (lines starting
%    with %). Write the prose under each comment block.
%  - Include it in the main file with % =============================================================================
%  CHAPTER: BACKGROUND
% =============================================================================
%  Master's thesis - Extension of "Balancing KPIs through Multi-Objective
%  Prescriptive Process Analytics" (Le, Buliga, Ronzani, Padella, de Leoni).
%
%  Thesis contribution: quantification of the predictive uncertainty
%  (aleatoric + epistemic) of the two CatBoost models, and use of "confidence
%  as a KPI" in the Pareto front, together with probability calibration and
%  diversified top-k selection.
%
%  HOW TO USE THIS FILE
%  -------------------
%  - This is only the skeleton of the chapter: every section/subsection is a
%    real LaTeX command; ALL the guidance text is in comments (lines starting
%    with %). Write the prose under each comment block.
%  - Include it in the main file with % =============================================================================
%  CHAPTER: BACKGROUND
% =============================================================================
%  Master's thesis - Extension of "Balancing KPIs through Multi-Objective
%  Prescriptive Process Analytics" (Le, Buliga, Ronzani, Padella, de Leoni).
%
%  Thesis contribution: quantification of the predictive uncertainty
%  (aleatoric + epistemic) of the two CatBoost models, and use of "confidence
%  as a KPI" in the Pareto front, together with probability calibration and
%  diversified top-k selection.
%
%  HOW TO USE THIS FILE
%  -------------------
%  - This is only the skeleton of the chapter: every section/subsection is a
%    real LaTeX command; ALL the guidance text is in comments (lines starting
%    with %). Write the prose under each comment block.
%  - Include it in the main file with % =============================================================================
%  CHAPTER: BACKGROUND
% =============================================================================
%  Master's thesis - Extension of "Balancing KPIs through Multi-Objective
%  Prescriptive Process Analytics" (Le, Buliga, Ronzani, Padella, de Leoni).
%
%  Thesis contribution: quantification of the predictive uncertainty
%  (aleatoric + epistemic) of the two CatBoost models, and use of "confidence
%  as a KPI" in the Pareto front, together with probability calibration and
%  diversified top-k selection.
%
%  HOW TO USE THIS FILE
%  -------------------
%  - This is only the skeleton of the chapter: every section/subsection is a
%    real LaTeX command; ALL the guidance text is in comments (lines starting
%    with %). Write the prose under each comment block.
%  - Include it in the main file with \input{background} (or \include{...}).
%  - Packages expected from the main file: booktabs (notation table), amsmath,
%    hyperref, a bibliography backend (biblatex or bibtex).
%  - Replace the \cite{...} keys with the ones in your .bib file.
%
%  GUIDING PRINCIPLES (note to self, not for the thesis)
%  ----------------------------------------------------
%  - FUNNEL structure: from the general (process mining) to the specific
%    (uncertainty in gradient boosting). Each section introduces ONLY what the
%    next one needs.
%  - Follow the system PIPELINE:
%       offline (log -> encoding -> predictive model)
%    -> online (feasibility -> multi-objective search -> recommendation)
%    -> evaluation (business process simulation).
%  - Boundary with the other chapters:
%     * BACKGROUND  = WHAT a concept is and WHY it exists (reusable literature,
%                     independent of your work).
%     * METHODOLOGY = how YOU combine those concepts (architecture, the
%                     uncertainty wrappers, confidence as a Pareto objective).
%     * EVALUATION / EXPERIMENTAL SETUP = the concrete VALUES (Optuna search
%                     space, n. of trials, window_size, tau/alpha thresholds,
%                     ...) and the result tables.
%  - HYPERPARAMETERS: in Background you explain ONLY what a hyperparameter is
%    and what each group controls. The chosen values and the search space go
%    in the Experimental Setup, with a pointer to \ref{sec:bg:gbm}.
%  - STATISTICAL tools do NOT get a section of their own: spread them across
%    the point where the system uses them (see "STATISTICAL TOOLS MAP" at the
%    end of the file).
%
%  SECTION DEPENDENCY ORDER
%  -----------------------
%    PM --+--> PPM --+--> GBM/CatBoost --> Uncertainty
%         |          +--> PrPA --> Multi-objective opt.
%         +----------------------------> Simulation
%    If a topic is ONLY yours   -> Methodology.
%    If it is reusable literature -> Background (even if little known: NSGA-II,
%    SGLB, temperature scaling, p-dispersion, conformal prediction).
% =============================================================================


\chapter{Background}
\label{chap:background}

This chapter introduces the concepts on which the rest of the thesis builds.
The problem it addresses lies at the intersection of several fields: process
mining provides the data model (event logs and traces); predictive and
prescriptive process monitoring turn a running trace into a recommendation;
gradient boosting is the predictive engine, and quantifying its uncertainty
is the core of this work; multi-objective optimization frames the trade-off
between conflicting KPIs; and business process simulation is the vehicle used
to evaluate the recommendations.

% -----------------------------------------------------------------------------
% (Chapter introduction above. Original planning notes kept for reference.)
% -----------------------------------------------------------------------------
% Goal: tell the reader WHAT they need to know and WHY, previewing the
% "offline -> online -> evaluation" thread and flagging that the central
% section is the one on uncertainty (the contribution). Close with a sentence
% on the chapter's structure.
%
% Outline:
%  - The problem tackled by the thesis lives at the intersection of process
%    mining, machine learning and multi-objective optimization.
%  - Introduce the domain first (event log, PPM, PrPA), then the predictive
%    tool (gradient boosting) and its uncertainty, finally the process
%    simulation used for evaluation.
%  - List of the sections.


% #############################################################################
\section{Process Mining: Basic Concepts}
\label{sec:bg:pm}
% #############################################################################
% Goal of the section: fix the minimal vocabulary (log, trace, event, prefix,
% activity, resource) on which everything else rests.
% References: \cite{vanderaalst2016} (Process Mining: Data Science in Action);
% \cite{adams2021} (concept drift, to motivate the BPIC2017 split).
% Do NOT put here: your transition-system cache implementation
% (utils/setup_cache.py) -> Methodology.
% Bridge to the next section: "Predictive models are built on top of these
% prefixes."

Process mining is a discipline that combines data science and process
management: it analyzes the data recorded by information systems to
reconstruct how business processes actually run~\cite{vanderaalst2016}. The
basic idea is simple. Every time a process is executed, for example a
reimbursement request, an order, a patient's hospital journey, the systems
that support it leave traces in the form of event logs: they record which
activity was performed, when, and for which specific case.

From this data, process mining techniques make it possible to automatically
reconstruct the process model as it is actually executed, to compare it with
an expected model to identify deviations, and to enrich it with additional
information useful for improving it~\cite{vanderaalst2016}. In practice, it
shows organizations the difference between how they think their processes work
and how they actually work.

\subsection{Event Log, Traces, Events}
\label{subsec:bg:pm:log}
The starting point of any process-mining analysis is the \emph{event log}. An
event log $\mathcal{L}$ collects the executions of a business process, called
\emph{traces}: each trace $\sigma = \langle e_1, e_2, \dots, e_n \rangle$ is a
chronologically ordered sequence of events produced by one instance (or case)
of the process, and $|\sigma|$ denotes its length. Every event
$e = (c, a, \mathbf{d})$ is uniquely tied to a trace through a case identifier
$c$, records the execution of one activity $a$, and carries a vector of
additional attributes $\mathbf{d}$ such as the start
and completion timestamps of the activity and the resource that performed
it~\cite{le2025balancing}.
$\mathcal{A}$ and $\mathcal{R}$ denote, respectively, the set of activity
labels and the set of resources that occur anywhere in the process;
$\mathcal{E}$ is the universe of possible events and $\mathcal{E}^{*}$ the
space of finite event sequences over it. Two projection functions,
$act : \mathcal{E} \to \mathcal{A}$ and $res : \mathcal{E} \to \mathcal{R}$,
extract the activity and the resource of an event. Because an activity can
occur more than once along the same trace (e.g.\ inside a loop), ``event''
and ``activity occurrence'' are the more precise terms, even though the
shorter ``activity'' is used informally throughout whenever no ambiguity
arises.

\subsection{Prefixes and Running Traces}
\label{subsec:bg:pm:prefix}
Historical event logs only contain completed traces, but a prescriptive
system observes a process while it is still unfolding. Given a complete trace
$\sigma$, its $k$-\emph{prefix} $\sigma_{(k)} = \langle e_1, \dots, e_k
\rangle$, with $k < |\sigma|$, is the sequence of its first $k$
events~\cite{le2025balancing}; the
remaining $\langle e_{k+1}, \dots, e_n \rangle$ is the suffix that has not
happened yet. A \emph{running}, or ongoing, execution is exactly such a
prefix, produced by a case that has not reached completion.

\subsection{State Abstraction and Transition Systems}
\label{subsec:bg:pm:ts}
% (Needed by Section \ref{sec:bg:prpa}.)
A running execution, as introduced above, is itself a prefix. A common way
to reason about it in process mining is a \emph{transition
system}~\cite{vanderaalst2016}, a
directed graph whose nodes are \emph{states} and whose edges are labeled
with activities: an edge from state $s$ to state $s'$ labeled $a$ means
that executing $a$ from $s$ leads to $s'$. States abstract prefixes
through a state-representation function $S$, mapping every prefix
$\sigma_{(k)}$ to some state $S(\sigma_{(k)})$; two prefixes that map to
the same state are treated as equivalent, even if they differ elsewhere.

Given a log of completed traces $\mathcal{L}$, the transition system is
mined by recording, for every state that occurs in the log, the
activities empirically observed to follow it. Formally, the set of
feasible next activities for a prefix $\sigma_{(k)}$
is~\cite{le2025balancing}
\begin{equation}
\mathcal{A}|_{\sigma_{(k)}} = \{a \in \mathcal{A} \mid \exists\, \sigma'
\in \mathcal{L}, \exists\, l < |\sigma'| : S(\tilde\sigma_{(l)}) =
S(\sigma_{(k)}), a = act(\sigma'_{(l+1)})\}.
\label{eq:bg:pm:feasible}
\end{equation}
That is, $a$ is feasible after $\sigma_{(k)}$ if some historical trace
reached the same state and was then followed by $a$. This is what lets
the transition system say something about a prefix it has never
literally seen, as long as its state was reached, at some point, by
another prefix in the log.

The choice of $S$ has direct consequences on the resulting transition
system. The most direct choice, using the prefix itself as the state
($S(\sigma_{(k)}) = \sigma_{(k)}$), offers little help: a state would
then correspond to one specific history, so nothing learned from one
case would transfer to another, and the number of states would grow
with every distinct prefix in the log. Some abstraction of the prefix
into a state is therefore needed, coarse enough to generalize across
similar histories, yet informative enough to keep the resulting
recommendation meaningful.

% -----------------------------------------------------------------------------
% Original planning notes for this section, kept for reference.
% -----------------------------------------------------------------------------
% subsec:bg:pm:log
% - Definition of event log L as a collection of traces (execution sequences
%   of a business process).
% - Trace sigma = <e1, ..., en>; |sigma| = number of events.
% - Event e = (c, a, d): case id c, activity a, data payload d (start/complete
%   timestamps, resource r, case attributes).
% - Projection functions act: E -> A and res: E -> R.
% - Sets A (activities) and R (resources) of the process.
% - Distinction event / activity / case; note that "event" and "activity
%   label" are often used interchangeably.
%
% subsec:bg:pm:prefix
% - k-prefix: sigma(k) = <e1, ..., ek> with k < |sigma|.
% - Running trace / ongoing execution: prefix of a trace that is not yet
%   complete; it is the input of the prescriptive problem.
% - Concept of "completion" of a prefix (its unknown suffix).
%
% subsec:bg:pm:ts
% - Why an abstraction is needed: avoid an over-deterministic transition
%   system and a state-space explosion.
% - Sequence-based state representation with window w:
%   S^w(sigma(k)) = last w activities of the prefix (or all of them if k < w).
% - Transition system mined from the log: maps abstract state -> activities
%   that can follow it.
% - Small example (reuse Fig. 2 of the paper or build one).


% #############################################################################
\section{Predictive Process Monitoring}
\label{sec:bg:ppm}
% #############################################################################
% Goal: how a property of the complete trace is predicted from a prefix;
% introduce the two targets used (outcome and remaining time) and the
% encoding.
% References: \cite{difrancescomarino2022} (PPM, Process Mining Handbook);
% \cite{marquezchamorro2018} (survey); \cite{teinemaa2018} (outcome-oriented);
% \cite{galanti2023}.
% Do NOT put here: your concrete sklearn pipeline, get_features() ->
% Methodology.
% Bridge: "PPM provides the estimate; Prescriptive Analytics uses it to
% decide what to do."

\subsection{The Predictive Process Monitoring Problem}
\label{subsec:bg:ppm:problem}
% - Definition: a function that, given a prefix, estimates a property of the
%   complete trace (outcome, remaining time, next activity, cost...).
% - Tasks relevant to the thesis:
%    * CASE OUTCOME: binary classification; defined by the occurrence of a key
%      activity (e.g. "O_Accepted" for BPIC2012/2017; "Network /
%      Back-Office Adjustment Requested" as a negative outcome for BAC).
%      The activity that defines the outcome is excluded from the
%      prescribable activities (fair evaluation).
%    * REMAINING TIME: regression on the time left until the case ends.
% - Notion of KPI function F(sigma, i): the KPI value after the first i events.

\subsection{Labeling}
\label{subsec:bg:ppm:labeling}
% - How the prefix training set is built: for every trace and every step k an
%   example (prefix, target) is generated.
% - Outcome label (boolean) and time label (continuous).

\subsection{Prefix Encoding}
\label{subsec:bg:ppm:encoding}
% - Why it is needed: tabular ML techniques require a fixed-length numeric
%   vector.
% - ACTIVITY-FREQUENCY ENCODING: a vector counting the occurrences of each
%   activity in the prefix. Other known options (last-state, aggregation,
%   index-based, embedding) to mention briefly for context.
% - Auxiliary attributes: prefix duration, current-event duration (end -
%   start, or fallback end - previous end), weekday of the last timestamp.
% - LABEL ENCODING of the candidate next activity and next resource
%   (NEXT_ACTIVITY, NEXT_RESOURCE): they are model inputs, not outputs.
% - Train/inference consistency: the same encoder is applied to the running
%   execution online.
% - One-hot vs label encoding: anticipate that the encoding choice affects
%   epistemic uncertainty (an unseen category becomes an all-zero row that
%   every sub-ensemble agrees on) -> picked up again in
%   \ref{sec:bg:uncertainty}.

\subsection{Normalization of the Time Target}
\label{subsec:bg:ppm:normalization}
% (BASIC STATISTICS.)
% - Problem: the remaining time has a strongly skewed (long-tail) distribution
%   and lives on a different scale than the binary outcome.
% - Three-step transformation (define it formally):
%    1. z-normalization (standardization): (x - mu) / s.
%    2. sigmoid: 1 / (1 + e^-z).
%    3. min-max scaling to [0, 1].
%   Result: the "sigmoid_mm" variable, nearly uniform on [0,1], which makes
%   the two KPIs comparable and makes even small improvements visible.
% - Invertibility of the transformation (needed to bring simulated times back
%   to seconds during evaluation): min-max^-1 -> logit -> de-standardization,
%   with edge clipping.
% - Recall: sample mean / standard deviation, the logistic function, monotone
%   transformations and order preservation.


% #############################################################################
\section{Prescriptive Process Analytics}
\label{sec:bg:prpa}
% #############################################################################
% Goal: introduce prescriptive process analytics (vs the other analytics
% types and vs predictive monitoring), give its formal problem statement, and
% explain what a KPI is. Three subsections: concept + examples; formal
% definition; KPIs.
% References: \cite{ibm_prescriptive} and \cite{wikipedia_prescriptive}
% (analytics taxonomy, technology behind prescriptive analytics -- consider
% replacing Wikipedia with a stronger source); \cite{kubrak2022} (Quo vadis?);
% \cite{le2025} (next step recommendation); \cite{le2025balancing} (the
% framework this thesis extends).
% Do NOT put here: the feasibility module internals, the search strategies
% (exhaustive vs genetic), the diverse top-k selection -> methodology chapter.
% Planning notes for that relocated material are kept as comments at the end
% of this section.
% Bridge: "When the KPIs are more than one and in conflict, 'optimal' must be
% redefined" -> Section \ref{sec:bg:moo}.

\subsection{From Prediction to Prescription}
\label{subsec:bg:prpa:pred2presc}
Data analytics is commonly described as a progression of four questions of
increasing sophistication and business value~\cite{ibm_prescriptive}:
\emph{descriptive} analytics asks \emph{what happened}, \emph{diagnostic}
analytics asks \emph{why it happened}, \emph{predictive} analytics asks
\emph{what is likely to happen next}, and \emph{prescriptive} analytics asks
\emph{what should be done next}. The four are complementary, each building on
the output of the previous one, but they differ in what they deliver: the
first three explain or forecast, whereas only the last one recommends a
course of action and makes the implications of each option
explicit~\cite{wikipedia_prescriptive}. Technically, prescriptive analytics
combines historical and contextual data with business rules and with
mathematical and computational models drawn from statistics, machine learning
and operations research, in order to reason about objectives, constraints,
uncertainty and trade-offs at
once~\cite{wikipedia_prescriptive,ibm_prescriptive}. It is applied, for
example, to recommend a treatment from a patient's characteristics, to set
prices and plan promotions from a demand forecast, to schedule maintenance
before a machine fails, or to score a transaction for fraud
risk~\cite{ibm_prescriptive}.

When this paradigm is applied to the executions of a business process
recorded in an event log, it is called \emph{Prescriptive Process Analytics}
(PrPA), or prescriptive process monitoring~\cite{kubrak2022}. PrPA supports
\emph{ongoing} process executions: instead of only forecasting how a running
case will end, it suggests interventions that steer it toward a desired
result, for instance preventing an unfavorable outcome or reducing the cost
or the duration of the case. The difference from predictive process
monitoring is one of intent. Predictive monitoring passively estimates a
property of the running case, such as its outcome, its remaining time or its
most likely next activity; prescriptive monitoring uses those estimates to
actively decide which action improves a target performance measure, taking
the dependencies between the different aspects of the process into account.

The kind of action prescribed varies across the literature. Early work
recommended the next activity to execute by exploiting the event data payload
and simple counterfactual reasoning~\cite{groger2014}. As prediction-based
methods spread, attention moved to flagging cases at risk of an undesired
outcome and to deciding \emph{whether} and \emph{when} to
intervene~\cite{teinemaa2018,fahrenkrogpetersen2022}. More recent approaches
recommend \emph{how} to proceed using causal inference~\cite{bozorgi2023},
reinforcement learning~\cite{branchi2022,metzger2020,shoush2022} or
interpretable, white-box predictors~\cite{padella2022}, and increasingly
prescribe not only the next activity but also the resource that should carry
it out --- the \emph{next step recommendation}~\cite{le2025} --- sometimes
accounting for how a resource's performance changes with
experience~\cite{padella2024}. Large language models have also started to be
used to make such recommendations easier to interpret~\cite{kubrak2024}.

What almost all of these approaches share is that they optimize a
\emph{single} performance measure at a time. The few that consider several
process aspects fold them into one fixed formula with a predetermined meaning
(e.g.\ monetary profit), which does not let a decision-maker examine
alternative trade-offs~\cite{le2025balancing}. Treating multiple, possibly
conflicting measures explicitly is the setting of this thesis; the tools for
it are introduced in Section~\ref{sec:bg:moo}.

\subsection{Formal Problem Definition}
\label{subsec:bg:prpa:formal}
A performance measure is formalized as a \emph{Key Performance Indicator}
(KPI) function $F : \mathcal{E}^{*} \times \mathbb{N} \to \mathbb{R}$ that,
given a trace $\sigma$ and a position $1 \le i \le |\sigma|$, returns the
value $F(\sigma, i)$ of that indicator for $\sigma$ after its first $i$
events~\cite{le2025balancing}. Evaluating a KPI at an intermediate position,
and not only on the completed trace, is what makes it usable while a case is
still running.

The prescription itself is a \emph{next step}: a pair
$(a, r) \in \mathcal{A} \times \mathcal{R}$ consisting of the activity to
perform next and the resource assigned to it. Given a running execution
$\sigma' = \langle e_1, \dots, e_k \rangle$, the prescriptive problem is to
choose the next event $e_{k+1}$ --- that is, the pair $(a, r)$ --- so that the
resulting completed trace
$\sigma_{\mathrm{pres}} = \langle e_1, \dots, e_k, e_{k+1}, \dots, e_n
\rangle$ attains the best possible KPI value. The candidate pairs are not
arbitrary: they are restricted to the activities that can actually follow the
current state of the case and to the resources able to perform them, a
feasibility constraint that is detailed, together with the rest of the
framework, in the methodology chapter. When a single KPI $F$ is considered,
``best'' simply means maximizing (or minimizing) $F$; when several KPIs
$F_1, \dots, F_n$ matter at once, they are aggregated into a single
\emph{utility function}
\begin{equation}
U_{\mathrm{KPI}}(\sigma, \boldsymbol{\lambda}, i)
  = \sum_{j=1}^{n} \lambda_j \, F_j(\sigma, i),
\label{eq:bg:prpa:utility}
\end{equation}
where $\boldsymbol{\lambda} = (\lambda_1, \dots, \lambda_n) \in [0, 1]^n$,
with $\sum_{j} \lambda_j = 1$, weights the relative importance of each
indicator according to the decision-maker's
priorities~\cite{le2025balancing}. The limitations of this linear
aggregation, and the alternative of not aggregating at all, are discussed in
Section~\ref{sec:bg:moo}.

\subsection{Key Performance Indicators}
\label{subsec:bg:prpa:kpi}
The objective of PrPA is stated informally as reaching a ``desired
outcome'', and a KPI is what turns that vague notion into a quantity that can
be optimized. Which quantity is appropriate depends entirely on the process
and on the organization's goals; recurring dimensions are the throughput time
of a case, its cost, some notion of quality or of regulatory compliance, and
the efficiency with which resources are used.

Two examples make the mapping concrete. In a loan-application process, a bank
that wants to grant more loans can take the \emph{acceptance rate} --- the
fraction of cases that reach a positive decision --- as a KPI to maximize,
whereas a bank that wants faster decisions can take the \emph{time to
decision} as a KPI to minimize. In a hospital, management may want to shorten
the \emph{length of stay} of patients in order to free capacity, while also
keeping the \emph{30-day readmission rate} low, so that discharging patients
earlier does not degrade the quality of care.

The two hospital KPIs illustrate a recurring difficulty: performance measures
often \emph{conflict}, so that improving one worsens another, and the right
balance between them is not fixed --- it depends on the domain, on the
regulatory context and on situational priorities. A prescription that is
optimal for one KPI can therefore be a poor choice overall. Making these
trade-offs explicit, instead of hiding them inside a single composite score,
is the problem addressed in Section~\ref{sec:bg:moo}.

% -----------------------------------------------------------------------------
% Planning notes -- material relocated to the methodology chapter, kept here
% for reference.
% -----------------------------------------------------------------------------
% Feasibility module:
% - Purpose: restrict the space to what is practically executable.
% - Activity filter: A|sigma(k) = activities that in the log follow the
%   abstract state of the prefix at least once (transition system, window w).
%   Fall back to progressively shorter suffixes if the exact prefix was never
%   seen.
% - Resource filter: R_a = resources that executed a at least once in the log.
% - Feasible prescription space: P_sigma(k) = {(a, r) : a in A|sigma(k),
%   r in R_a}.
% - Exclusion of "forbidden" activities (those defining the outcome).
%
% Search-space exploration strategies:
% - Exhaustive search: evaluates all feasible pairs; optimal w.r.t. the model
%   but costly when |P_sigma(k)| is large (many resources).
% - Heuristic / metaheuristic (genetic) methods: explore a subset.
%   (Algorithm detail belongs to Section \ref{sec:bg:moo} / methodology.)
%
% Extended single-KPI survey (if a separate Related Work chapter is wanted):
% - Resource allocation for global KPI improvement \cite{meneghello2024}.
% - de Leoni et al. on guidance for process continuation \cite{deleoni2020}.


% #############################################################################
\section{Conflicting KPIs and Multi-Objective Optimization}
\label{sec:bg:moo}
% #############################################################################
% Goal: the tools to understand why a Pareto front is used, how it is explored
% (enumeration, NSGA-II) and how a diversified subset of solutions is chosen.
% References: \cite{marler2004} (multi-objective survey); \cite{deb2002}
% (NSGA-II); \cite{mitchell1998} (genetic algorithms); \cite{mothilal2020}
% (DiCE); \cite{buliga2025}; p-dispersion / facility-location literature.
% Do NOT put here: the layout of the recommendations_*.csv files and the
% _objectives.csv sidecars -> Methodology / Evaluation.
% Bridge: "The quality of these decisions depends entirely on the quality -
% and the reliability - of the model that estimates the three objectives."

So far, most prescriptive-monitoring approaches optimize a single KPI. In
practice this is rarely enough: a business process is judged on several
performance measures at once, these measures are usually interdependent, and
improving one of them often comes at the expense of another. What is needed is
a way to reason about the whole set of measures together and to choose a
deliberate trade-off among them.

\subsection{Conflicting KPIs}
\label{subsec:bg:moo:devil}
The tension is easiest to see on an example. A hospital may want to shorten
the length of stay of its patients, in order to free capacity, and at the
same time keep the 30-day readmission rate low; discharging patients sooner
helps the first measure but tends to hurt the second, so no policy is best on
both counts and management has to settle on an acceptable compromise. The
same situation recurs across domains, and it is not limited to two measures:
the classical \emph{Devil's Quadrangle} already identifies four dimensions of
process performance --- time, cost, quality and flexibility --- that
typically cannot all be improved at once~\cite{dumas2018}.

\subsection{Scalarization and Its Limitations}
\label{subsec:bg:moo:scalarization}
One way to keep the convenience of a single optimum is to collapse the
several KPIs into one number. This is what the utility function of
Equation~\eqref{eq:bg:prpa:utility} does: a weighted sum of the individual
indicators, with weights $\boldsymbol{\lambda}$ chosen by the
decision-maker~\cite{marler2004}. Scalarization of this kind is simple and it
always yields one well-defined solution, but it has three drawbacks. It
forces the weights to be fixed \emph{before} any solution is seen, even
though the trade-off they encode is exactly what the decision-maker would
like to inspect; it returns a single point for each weight vector, hiding the
shape of the trade-off; and it requires the KPIs to be expressed on
comparable scales, which is why the time target is normalized as described in
Section~\ref{subsec:bg:ppm:normalization}.

\subsection{Pareto Optimality}
\label{subsec:bg:moo:pareto}
When the objectives genuinely conflict there is no single solution that is
best on all of them, and the problem instead has a whole set of equally
defensible compromises. This is captured by the notion of \emph{Pareto
optimality}~\cite{pareto1906,miettinen1999}. Consider a minimization problem
with $m$ objective functions $f_1, \dots, f_m$ and a decision vector $x$
ranging over a feasible set $X$. A solution $x_1$ is said to \emph{dominate}
a solution $x_2$, written $x_1 \prec x_2$, when it is no worse on every
objective and strictly better on at least one:
\begin{equation}
x_1 \prec x_2
\iff
\bigl( f_i(x_1) \le f_i(x_2)\ \ \forall\, i \in \{1, \dots, m\} \bigr)
\ \wedge\
\bigl( \exists\, j : f_j(x_1) < f_j(x_2) \bigr).
\label{eq:bg:moo:dominance}
\end{equation}
A solution is \emph{Pareto-optimal} --- equivalently \emph{efficient} or
\emph{non-dominated} --- if no feasible solution dominates it. The image of
the Pareto-optimal set in objective space is the \emph{Pareto front}
\begin{equation}
P^{*} = \bigl\{\, f(x) \;\big|\; x \in X,\ \nexists\, x' \in X :
f(x') \prec f(x) \,\bigr\}.
\label{eq:bg:moo:front}
\end{equation}
A point lies on the front simply because nothing dominates it; it need not
itself dominate any other point. The best and the worst value that each
objective attains over $P^{*}$ define the \emph{ideal} and the \emph{nadir}
point, which are convenient references for normalizing the objectives and for
measuring how close a given solution is to the (generally unattainable)
ideal. When two entire fronts have to be compared rather than two individual
points, a common scalar summary is the \emph{hypervolume}: the volume of
objective space that a front dominates with respect to a fixed reference
point~\cite{deb2001}.

Compared with scalarization, computing the Pareto front defers the trade-off
decision. The optimizer returns all the non-dominated compromises, and the
selection of one of them is made afterwards, with the cost of each
alternative in full view~\cite{honegumi}. In this thesis the front is built
on three objectives: the probability of a positive outcome, to be maximized;
the predicted remaining time, to be minimized; and the predictive standard
deviation of that time estimate, also to be minimized. The third objective
treats the model's own \emph{confidence} as a KPI, a choice motivated in
Section~\ref{sec:bg:uncertainty}.

\subsection{Front-Exploration Algorithms}
\label{subsec:bg:moo:algorithms}
% - Exact enumeration + front extraction: exact on small discrete spaces.
% - GENETIC ALGORITHMS: population, fitness, selection, crossover, mutation,
%   elitism, stopping criterion.
%    * The original paper's "custom fitness" approach: validity (val) +
%      alpha * proximity (prox), inspired by \cite{buliga2025}; DiCE as the
%      library \cite{mothilal2020}.
%    * NSGA-II \cite{deb2002}: non-dominated sorting into fronts, crowding
%      distance for diversity, tournament selection, elitism.
%      Implemented with pymoo (an integer variable indexing the (a,r) pair,
%      SBX + polynomial mutation + rounding repair), currently "parked"
%      because it is slower on small candidate sets: explain it anyway, it is
%      a methodological contribution.
% - Enumeration vs metaheuristic comparison: quality vs cost; robustness to
%   predictor errors on rare pairs (the paper's hypothesis: proximity favors
%   frequent pairs, which are estimated better).

\subsection{Selecting a Representative Subset of the Front}
\label{subsec:bg:moo:dispersion}
For a discrete problem the Pareto front can contain a large number of points,
some of them very close to one another. A recommender that has to present
alternatives to a user is often better served by a handful of clearly
distinct options than by the entire front. Choosing the $k$ most
representative points can be framed as a \emph{max-min dispersion}, or
\emph{$p$-dispersion}, problem: given $n$ candidate points with pairwise
distances $d_{ij}$, select a subset $S$ of size $k$ that maximizes the
smallest distance between any two selected points,
\begin{equation}
\max_{\substack{S \subseteq \{1, \dots, n\} \\ |S| = k}}\ \
\min_{\substack{i, j \in S \\ i \neq j}}\ d_{ij}.
\label{eq:bg:moo:dispersion}
\end{equation}
The problem was first studied in facility location, where the aim is to place
$p$ facilities so that the closest pair is as far apart as
possible~\cite{kuby1987}, and it has since been used more generally to
extract a maximally spread-out subset from a set of
candidates~\cite{maliszewski2012}. Because it is stated purely in terms of a
matrix of pairwise distances, the distance metric is a modeling choice; here
the natural one is the Euclidean distance between points in the normalized
objective space. The problem admits an exact mixed-integer programming
formulation, for which off-the-shelf solvers are
available~\cite{feng2022spopt}; the formulation actually used, together with
the rule that picks a single final recommendation, is given in the
methodology chapter.

Every objective on the front is produced by a predictive model, so the value
of the whole procedure rests on how accurate --- and how well calibrated ---
that model is. The next two sections describe it.


% #############################################################################
\section{Gradient Boosting and CatBoost}
\label{sec:bg:gbm}
% #############################################################################
% Goal: the model class used for the two predictors and what hyperparameters
% are / what role they play, plus their optimization. NO concrete values here.
% References: \cite{friedman2001} (gradient boosting); \cite{prokhorenkova2018}
% (CatBoost); \cite{akiba2019} (Optuna); \cite{bergstra2011} (TPE);
% \cite{galanti2023}.
% Do NOT put here: the search_spaces of 2_training_predictive_model.py, the
% number of trials, the best_hyperparams_*.json -> Experimental Setup (with a
% pointer to \ref{sec:bg:gbm}).
% Bridge: "A point model does not say how sure it is of what it predicts:
% uncertainty must be quantified."

\subsection{Tree Ensembles and Boosting}
\label{subsec:bg:gbm:boosting}
% - Decision / regression tree: splits, leaves, depth.
% - Bagging vs boosting (general idea).
% - Gradient boosting: sequential addition of trees fitting the negative
%   gradient of the loss; learning rate as shrinkage; number of trees
%   (iterations) and overfitting.

\subsection{CatBoost}
\label{subsec:bg:gbm:catboost}
% - Motivation for the choice: high predictive quality with contained
%   training time; native handling of categorical features; used in prior PPM
%   work \cite{galanti2023}.
% - Distinctive elements:
%    * ordered boosting / ordered target statistics (mitigates target leakage
%      / prediction shift);
%    * oblivious (symmetric) trees: same split at every level -> balanced
%      trees, fast inference, regularizing effect;
%    * growth alternatives (Depthwise, Lossguide) and min_data_in_leaf.
% - Losses used in the thesis (as CONCEPTS):
%    * Logloss for outcome classification;
%    * RMSE for regression; and the RMSEWithUncertainty variant, which
%      predicts mean AND variance (bridge to \ref{sec:bg:uncertainty}).
% - scale_pos_weight / class weights for imbalance (the positive-case
%   percentages differ a lot across logs: e.g. BAC ~82.5%, BPIC2012 ~47.9%).

\subsection{Hyperparameters: What They Are and What They Control}
\label{subsec:bg:gbm:hyperparams}
% - Definition: a parameter NOT learned from the data, fixed before training,
%   governing model capacity and the optimization process.
% - Grouped by ROLE (explain the role, not the value):
%    * Capacity / complexity: depth, grow_policy, min_data_in_leaf.
%    * Learning step: learning_rate, iterations.
%    * Regularization: l2_leaf_reg.
%    * Stochasticity / sub-sampling: bootstrap_type (Bayesian/Bernoulli/MVS),
%      bagging_temperature, subsample, colsample_bylevel.
%    * Imbalance: scale_pos_weight.
% - The bias-variance trade-off is governed by these groups.

\subsection{Hyperparameter Optimization}
\label{subsec:bg:gbm:hpo}
% - Why not by hand: large space, non-linear interactions.
% - Validation: hold-out, k-fold cross-validation, and - for processes - a
%   TEMPORAL split (avoid leakage from the future). Note the internal
%   "running" validation (the most recent training cases as validation).
% - Early stopping: stop adding trees when the validation metric stops
%   improving.
% - Search: grid / random search vs BAYESIAN OPTIMIZATION.
%    * Optuna: TPE (Tree-structured Parzen Estimator) as the surrogate;
%      define-by-run search space.
%    * Pruning (MedianPruner): interrupt unpromising trials.
% - Selection of the winning trial on ONE consistent metric (Logloss / RMSE)
%   and a final refit on 100% of the training data with the chosen tree count.


% #############################################################################
\section{Predictive Uncertainty Quantification}
\label{sec:bg:uncertainty}
% #############################################################################
% Goal: this is the THEORETICAL CORE of the contribution. Taxonomy of
% uncertainty, its estimation in ensembles and in gradient boosting,
% calibration. It must be the most developed section of the chapter.
% References: \cite{malinin2021} (Uncertainty in Gradient Boosting via
% Ensembles, ICLR 2021); \cite{hullermeier2021} (aleatoric vs epistemic,
% survey); \cite{lakshminarayanan2017} (deep ensembles); \cite{guo2017}
% (temperature scaling); \cite{kuleshov2018} (calibrated regression);
% \cite{angelopoulos2023} (conformal prediction); \cite{gal2016} (MC-dropout).
% Do NOT put here: the UncertaintyRegressor / UncertaintyClassifier classes,
% how the wrappers plug into the Pipeline, how out["std"] becomes objective #3
% -> Methodology. The obtained ECE/coverage/T/s values -> Evaluation.
% Bridge: "The Pareto front and the recommendations are evaluated with a
% process simulation, because a real A/B test is not possible."

\subsection{Taxonomy of Uncertainty}
\label{subsec:bg:uncertainty:taxonomy}
% - ALEATORIC uncertainty (data uncertainty): intrinsic / irreducible noise in
%   the data-generating process; does not shrink with more data.
%    * homoscedastic vs heteroscedastic.
% - EPISTEMIC uncertainty (knowledge / model uncertainty): the model's
%   ignorance due to lack of data / coverage; reducible with more data; high
%   on rare or out-of-distribution inputs.
% - TOTAL predictive uncertainty: the combination of the two.
% - Why distinguish them in the thesis: epistemic flags out-of-domain inputs /
%   rarely observed (a,r) pairs; aleatoric says how noisy the estimate is even
%   in the ideal case.
% - RESULT PREVIEW: on one-hot tabular data the regressor's epistemic part is
%   very small (~0.1% of the variance) and does NOT rise on never-seen inputs
%   (an unknown category -> an all-zero row every sub-ensemble agrees on). It
%   is reported anyway: the full decomposition is a result in itself
%   ("epistemic is negligible here, and why").

\subsection{Ensemble Methods for Uncertainty}
\label{subsec:bg:uncertainty:ensembles}
% - Idea: train/extract several models; the DISAGREEMENT among members
%   estimates the epistemic part, the AVERAGE member variance estimates the
%   aleatoric part.
% - Deep ensembles, MC-dropout, bootstrap ensembles: cite as context.
% - LAW OF TOTAL VARIANCE (regression):
%     Var[y] = E_m[ Var(y | m) ]  +  Var_m[ E(y | m) ]
%              \____ aleatoric ___/    \___ epistemic ___/
% - ENTROPY decomposition (classification):
%     H[ E_m p_m ]            = total uncertainty
%     E_m[ H(p_m) ]           = aleatoric uncertainty (expected entropy)
%     MI = total - aleatoric  = epistemic uncertainty (mutual information, BALD)
%   All in nats.

\subsection{Uncertainty in Gradient Boosting}
\label{subsec:bg:uncertainty:gbm}
% (\cite{malinin2021}.)
% - SGLB - Stochastic Gradient Langevin Boosting: noise injected into the
%   updates -> the model (approximately) samples from a posterior over the
%   parameters. In CatBoost: posterior_sampling=True.
% - VIRTUAL ENSEMBLES: from a single SGLB model several "sub-models" are
%   obtained using growing prefixes of the tree sequence -> an ensemble at the
%   cost of a single training run. Parameter: virtual_ensembles_count.
% - Regression with the RMSEWithUncertainty loss: each sub-model predicts
%   (mu_m, sigma_m^2); via the law of total variance:
%     data_std      = sqrt( mean_m sigma_m^2 )     (aleatoric)
%     knowledge_std = sqrt( var_m mu_m )           (epistemic)
%     total_std     = sqrt( data_var + knowledge_var )
% - Classification: prediction_type="TotalUncertainty" returns the expected
%   entropy (aleatoric) and the total entropy; the mutual information
%   (epistemic) is the difference.

\subsection{Calibration}
\label{subsec:bg:uncertainty:calibration}
% - What a CALIBRATED (probabilistic) model is: among the events predicted
%   with probability p, roughly a fraction p actually occur.
% - Diagnostics:
%    * reliability diagram;
%    * ECE - Expected Calibration Error: |confidence - accuracy| averaged over
%      probability bins;
%    * for regression: interval COVERAGE (fraction of |y - mu| within
%      +-1 sigma / +-2 sigma; target ~68% / ~95% for a Gaussian) and
%      SHARPNESS (typical / median interval width).
% - POST-HOC correction methods:
%    * TEMPERATURE SCALING \cite{guo2017}: a single scalar T,
%      p_cal = sigmoid( logit(p) / T ); T > 1 "softens" over-confident
%      probabilities. Monotone -> it does not reorder candidates by
%      probability, but it changes the magnitudes feeding the normalization
%      and the front distances.
%    * SIGMA RECALIBRATION for regression: a scalar factor s that rescales
%      every standard-deviation component so the empirical coverage matches
%      the nominal one; estimated on a held-out slice with a robust
%      median-of-standardized-residual estimator (avoids blow-ups when
%      CatBoost predicts a near-zero variance).
% - CONFORMAL PREDICTION as the RIGOROUS alternative \cite{angelopoulos2023}:
%   split-conformal on a dedicated holdout -> intervals with a distribution-
%   free marginal coverage guarantee. Cite it as "the rigorous alternative"
%   and motivate why the thesis used post-hoc recalibration instead (the
%   models are trained once, cost/benefit trade-off).

\subsection{Supporting Statistical Tools}
\label{subsec:bg:uncertainty:stats}
% (Compact recap.)
% - Variance and standard deviation as dispersion measures.
% - Shannon entropy of a Bernoulli; entropy as an uncertainty measure.
% - Mutual information / information gain.
% - The normal distribution, the 68-95-99.7 rule, the standardized residual,
%   the 0.6745 quantile of the normal (median absolute value).


% #############################################################################
\section{Business Process Simulation}
\label{sec:bg:simulation}
% #############################################################################
% Goal: what process simulation is, how a simulation model is obtained from
% the data, and with which metrics its fidelity is assessed.
% References: \cite{vinci2025prosit} (ProSiT); \cite{demoor2025} (SimBank);
% \cite{chapelacampa2024} (extraneous delays); \cite{vanderaalst2016} (Petri
% nets / soundness); \cite{vinci2025online} (online discovery).
% Do NOT put here: the prosit_simulation_results/... paths, the
% build_recommender_df code, the specific case ids -> Methodology / Evaluation.

\subsection{Why Simulation}
\label{subsec:bg:simulation:why}
% - Impossibility of an A/B test: one cannot intervene on a real process just
%   to evaluate.
% - COUNTERFACTUAL evaluation: simulate "what would have happened" following /
%   not following the recommendation.
% - Standard in PrPA for evaluating recommendations \cite{demoor2025}.

\subsection{Petri Nets and Process Models}
\label{subsec:bg:simulation:petri}
% - Petri net: places, transitions, arcs, marking, firing of a transition.
% - Workflow net and soundness (every case can always complete).
% - Initial / final marking; relation to the start/end of a case.
% - Replay and ALIGNMENT of a log prefix against the model: reconstruct the
%   marking reached by a historical prefix (log move, model move,
%   model-inserted activities).

\subsection{Stochastic Discrete-Event Simulation}
\label{subsec:bg:simulation:des}
% - Components of a process simulation model:
%    * case ARRIVAL process;
%    * transition branching weights / probabilities;
%    * activity EXECUTION times;
%    * WAITING times / queues;
%    * resource CALENDARS and availability;
%    * (extraneous delays \cite{chapelacampa2024}, as context).
% - Repeating the runs to mitigate stochasticity (e.g. 10 runs per scenario).

\subsection{Discovery of the Simulation Model}
\label{subsec:bg:simulation:discovery}
% - The simulation parameters are LEARNED from the historical log (resource
%   discovery, transition weights, time distributions, calendars).
% - METHODOLOGICAL NOTE: the simulator parameters are discovered from the FULL
%   log, not from the training split.
% - The Petri-net structure is a separate input (control-flow discovery),
%   distinct from the parameter discovery.
% - Online discovery for evolving processes \cite{vinci2025online} - context.

\subsection{ProSiT}
\label{subsec:bg:simulation:prosit}
% - The simulator used \cite{vinci2025prosit}: interactive and transparent
%   simulations.
% - What it returns: simulated logs; diagnostics (unreachable recommendations,
%   truncated "runaway" cases, non-fitting prefixes, model-inserted
%   activities).

\subsection{Assessing the Simulator's Fidelity}
\label{subsec:bg:simulation:fidelity}
% (STATISTICS.)
% - Real vs simulated distribution comparison:
%    * case-duration distributions (histograms / KDE, median);
%    * rate of positive-outcome cases, real vs simulated.
% - Distribution of the simulator's residuals.
% - (Optional) distribution-distance metrics: EMD / Wasserstein, KS; or a
%   comparison of summary statistics.

\subsection{KPI Improvement Metrics}
\label{subsec:bg:simulation:kpimetrics}
% (Definition; the numbers are in Evaluation.)
% - Delta_CO = average over prefixes of ( F_CO(sigma_pres) - F_CO(sigma_GT) ).
% - Delta_RT = average relative improvement of the remaining time w.r.t. the
%   ground truth.
% - Baseline / Ground Truth: no recommendation applied.
% - (Optional) confidence intervals on the deltas via bootstrap; 95%
%   confidence bands in the plots as lambda varies.


% #############################################################################
\section{Notation}
\label{sec:bg:notation}
% #############################################################################
% OPTIONAL but much appreciated by advisors. A half-page table of symbols,
% consistent with those used in the rest of the thesis.
% Uncomment and complete the table below (requires \usepackage{booktabs}).
%
% \begin{table}[h]
%   \centering
%   \begin{tabular}{ll}
%     \toprule
%     Symbol & Meaning \\
%     \midrule
%     $\mathcal{L}$              & event log \\
%     $\sigma,\ \sigma_{(k)}$    & trace, $k$-prefix \\
%     $\mathcal{A},\ \mathcal{R}$ & sets of activities and resources \\
%     $act,\ res$                & projection functions \\
%     $e = (c, a, \mathbf{d})$   & event (case id, activity, payload) \\
%     $S^{w}(\sigma_{(k)})$      & state abstraction with window $w$ \\
%     $\mathcal{P}_{\sigma_{(k)}}$ & feasible prescription space \\
%     $F_{CO},\ F_{RT}$          & KPIs: case outcome, remaining time (norm.) \\
%     $U_{KPI}$                  & utility function / combined KPI \\
%     $\lambda$                  & KPI trade-off weight \\
%     $\phi$                     & predictive model (regressor) \\
%     $p,\ p_{cal}$              & raw / calibrated outcome probability \\
%     $T$                        & calibration temperature \cite{guo2017} \\
%     $s$                        & std recalibration factor \\
%     $\sigma_{data}$            & aleatoric uncertainty (std) \\
%     $\sigma_{know}$            & epistemic uncertainty (std) \\
%     $\sigma_{tot}$             & total predictive uncertainty \\
%     $\mathrm{MI}$              & mutual information (epistemic, classif.) \\
%     $k$                        & n. of diversified recommendations per case \\
%     $\tau,\ \alpha$            & validity threshold, proximity weight \\
%     \bottomrule
%   \end{tabular}
%   \caption{Notation used in the thesis.}
%   \label{tab:notation}
% \end{table}


% =============================================================================
%  STATISTICAL TOOLS MAP -> WHERE THEY GO  (note to self, not for the thesis)
% =============================================================================
%   z-score, sigmoid, min-max, invertibility             -> \ref{subsec:bg:ppm:normalization}
%   Bayesian optimization / TPE, CV, temporal split,
%     early stopping, pruning                            -> \ref{subsec:bg:gbm:hpo}
%   variance / std, Shannon entropy, mutual information,
%     law of total variance                              -> \ref{subsec:bg:uncertainty:ensembles} / \ref{subsec:bg:uncertainty:stats}
%   ECE, reliability diagram, temperature scaling,
%     coverage, sharpness, conformal prediction          -> \ref{subsec:bg:uncertainty:calibration}
%   68-95-99.7 rule, standardized residual,
%     0.6745 quantile                                    -> \ref{subsec:bg:uncertainty:stats}
%   distribution comparison (KDE, median, EMD/KS),
%     residuals                                          -> \ref{subsec:bg:simulation:fidelity}
%   bootstrap / 95% confidence bands                     -> \ref{subsec:bg:simulation:kpimetrics} and Evaluation
%
%   HYPERPARAMETERS appear in Background ONLY as concepts
%   (\ref{subsec:bg:gbm:hyperparams}). Chosen values, search space, number of
%   Optuna trials and best-params tables -> Experimental Setup, with the
%   sentence "the concepts are introduced in Section~\ref{sec:bg:gbm}".
%
% =============================================================================
%  FINAL CHECKLIST BEFORE HANDING IN THE CHAPTER
% =============================================================================
%  [ ] Every section ends with a "bridge" to the next one.
%  [ ] No concrete experimental value in Background (concepts only).
%  [ ] No description of YOUR classes/files/functions (-> Methodology).
%  [ ] Every concept introduced is later used in Methodology or Evaluation
%      (no "orphan" background).
%  [ ] \ref{sec:bg:uncertainty} is the most developed section: it is the
%      contribution.
%  [ ] The aleatoric/epistemic taxonomy is explained BEFORE SGLB / virtual
%      ensembles.
%  [ ] Notation table consistent with the symbols in the rest of the thesis.
%  [ ] Every subsection has at least one solid bibliographic reference.
%  [ ] Agreed with the advisors whether Related Work stays here
%      (\ref{subsec:bg:prpa:sota}) or is a separate chapter.
% =============================================================================
 (or \include{...}).
%  - Packages expected from the main file: booktabs (notation table), amsmath,
%    hyperref, a bibliography backend (biblatex or bibtex).
%  - Replace the \cite{...} keys with the ones in your .bib file.
%
%  GUIDING PRINCIPLES (note to self, not for the thesis)
%  ----------------------------------------------------
%  - FUNNEL structure: from the general (process mining) to the specific
%    (uncertainty in gradient boosting). Each section introduces ONLY what the
%    next one needs.
%  - Follow the system PIPELINE:
%       offline (log -> encoding -> predictive model)
%    -> online (feasibility -> multi-objective search -> recommendation)
%    -> evaluation (business process simulation).
%  - Boundary with the other chapters:
%     * BACKGROUND  = WHAT a concept is and WHY it exists (reusable literature,
%                     independent of your work).
%     * METHODOLOGY = how YOU combine those concepts (architecture, the
%                     uncertainty wrappers, confidence as a Pareto objective).
%     * EVALUATION / EXPERIMENTAL SETUP = the concrete VALUES (Optuna search
%                     space, n. of trials, window_size, tau/alpha thresholds,
%                     ...) and the result tables.
%  - HYPERPARAMETERS: in Background you explain ONLY what a hyperparameter is
%    and what each group controls. The chosen values and the search space go
%    in the Experimental Setup, with a pointer to \ref{sec:bg:gbm}.
%  - STATISTICAL tools do NOT get a section of their own: spread them across
%    the point where the system uses them (see "STATISTICAL TOOLS MAP" at the
%    end of the file).
%
%  SECTION DEPENDENCY ORDER
%  -----------------------
%    PM --+--> PPM --+--> GBM/CatBoost --> Uncertainty
%         |          +--> PrPA --> Multi-objective opt.
%         +----------------------------> Simulation
%    If a topic is ONLY yours   -> Methodology.
%    If it is reusable literature -> Background (even if little known: NSGA-II,
%    SGLB, temperature scaling, p-dispersion, conformal prediction).
% =============================================================================


\chapter{Background}
\label{chap:background}

This chapter introduces the concepts on which the rest of the thesis builds.
The problem it addresses lies at the intersection of several fields: process
mining provides the data model (event logs and traces); predictive and
prescriptive process monitoring turn a running trace into a recommendation;
gradient boosting is the predictive engine, and quantifying its uncertainty
is the core of this work; multi-objective optimization frames the trade-off
between conflicting KPIs; and business process simulation is the vehicle used
to evaluate the recommendations.

% -----------------------------------------------------------------------------
% (Chapter introduction above. Original planning notes kept for reference.)
% -----------------------------------------------------------------------------
% Goal: tell the reader WHAT they need to know and WHY, previewing the
% "offline -> online -> evaluation" thread and flagging that the central
% section is the one on uncertainty (the contribution). Close with a sentence
% on the chapter's structure.
%
% Outline:
%  - The problem tackled by the thesis lives at the intersection of process
%    mining, machine learning and multi-objective optimization.
%  - Introduce the domain first (event log, PPM, PrPA), then the predictive
%    tool (gradient boosting) and its uncertainty, finally the process
%    simulation used for evaluation.
%  - List of the sections.


% #############################################################################
\section{Process Mining: Basic Concepts}
\label{sec:bg:pm}
% #############################################################################
% Goal of the section: fix the minimal vocabulary (log, trace, event, prefix,
% activity, resource) on which everything else rests.
% References: \cite{vanderaalst2016} (Process Mining: Data Science in Action);
% \cite{adams2021} (concept drift, to motivate the BPIC2017 split).
% Do NOT put here: your transition-system cache implementation
% (utils/setup_cache.py) -> Methodology.
% Bridge to the next section: "Predictive models are built on top of these
% prefixes."

Process mining is a discipline that combines data science and process
management: it analyzes the data recorded by information systems to
reconstruct how business processes actually run~\cite{vanderaalst2016}. The
basic idea is simple. Every time a process is executed, for example a
reimbursement request, an order, a patient's hospital journey, the systems
that support it leave traces in the form of event logs: they record which
activity was performed, when, and for which specific case.

From this data, process mining techniques make it possible to automatically
reconstruct the process model as it is actually executed, to compare it with
an expected model to identify deviations, and to enrich it with additional
information useful for improving it~\cite{vanderaalst2016}. In practice, it
shows organizations the difference between how they think their processes work
and how they actually work.

\subsection{Event Log, Traces, Events}
\label{subsec:bg:pm:log}
The starting point of any process-mining analysis is the \emph{event log}. An
event log $\mathcal{L}$ collects the executions of a business process, called
\emph{traces}: each trace $\sigma = \langle e_1, e_2, \dots, e_n \rangle$ is a
chronologically ordered sequence of events produced by one instance (or case)
of the process, and $|\sigma|$ denotes its length. Every event
$e = (c, a, \mathbf{d})$ is uniquely tied to a trace through a case identifier
$c$, records the execution of one activity $a$, and carries a vector of
additional attributes $\mathbf{d}$ such as the start
and completion timestamps of the activity and the resource that performed
it~\cite{le2025balancing}.
$\mathcal{A}$ and $\mathcal{R}$ denote, respectively, the set of activity
labels and the set of resources that occur anywhere in the process;
$\mathcal{E}$ is the universe of possible events and $\mathcal{E}^{*}$ the
space of finite event sequences over it. Two projection functions,
$act : \mathcal{E} \to \mathcal{A}$ and $res : \mathcal{E} \to \mathcal{R}$,
extract the activity and the resource of an event. Because an activity can
occur more than once along the same trace (e.g.\ inside a loop), ``event''
and ``activity occurrence'' are the more precise terms, even though the
shorter ``activity'' is used informally throughout whenever no ambiguity
arises.

\subsection{Prefixes and Running Traces}
\label{subsec:bg:pm:prefix}
Historical event logs only contain completed traces, but a prescriptive
system observes a process while it is still unfolding. Given a complete trace
$\sigma$, its $k$-\emph{prefix} $\sigma_{(k)} = \langle e_1, \dots, e_k
\rangle$, with $k < |\sigma|$, is the sequence of its first $k$
events~\cite{le2025balancing}; the
remaining $\langle e_{k+1}, \dots, e_n \rangle$ is the suffix that has not
happened yet. A \emph{running}, or ongoing, execution is exactly such a
prefix, produced by a case that has not reached completion.

\subsection{State Abstraction and Transition Systems}
\label{subsec:bg:pm:ts}
% (Needed by Section \ref{sec:bg:prpa}.)
A running execution, as introduced above, is itself a prefix. A common way
to reason about it in process mining is a \emph{transition
system}~\cite{vanderaalst2016}, a
directed graph whose nodes are \emph{states} and whose edges are labeled
with activities: an edge from state $s$ to state $s'$ labeled $a$ means
that executing $a$ from $s$ leads to $s'$. States abstract prefixes
through a state-representation function $S$, mapping every prefix
$\sigma_{(k)}$ to some state $S(\sigma_{(k)})$; two prefixes that map to
the same state are treated as equivalent, even if they differ elsewhere.

Given a log of completed traces $\mathcal{L}$, the transition system is
mined by recording, for every state that occurs in the log, the
activities empirically observed to follow it. Formally, the set of
feasible next activities for a prefix $\sigma_{(k)}$
is~\cite{le2025balancing}
\begin{equation}
\mathcal{A}|_{\sigma_{(k)}} = \{a \in \mathcal{A} \mid \exists\, \sigma'
\in \mathcal{L}, \exists\, l < |\sigma'| : S(\tilde\sigma_{(l)}) =
S(\sigma_{(k)}), a = act(\sigma'_{(l+1)})\}.
\label{eq:bg:pm:feasible}
\end{equation}
That is, $a$ is feasible after $\sigma_{(k)}$ if some historical trace
reached the same state and was then followed by $a$. This is what lets
the transition system say something about a prefix it has never
literally seen, as long as its state was reached, at some point, by
another prefix in the log.

The choice of $S$ has direct consequences on the resulting transition
system. The most direct choice, using the prefix itself as the state
($S(\sigma_{(k)}) = \sigma_{(k)}$), offers little help: a state would
then correspond to one specific history, so nothing learned from one
case would transfer to another, and the number of states would grow
with every distinct prefix in the log. Some abstraction of the prefix
into a state is therefore needed, coarse enough to generalize across
similar histories, yet informative enough to keep the resulting
recommendation meaningful.

% -----------------------------------------------------------------------------
% Original planning notes for this section, kept for reference.
% -----------------------------------------------------------------------------
% subsec:bg:pm:log
% - Definition of event log L as a collection of traces (execution sequences
%   of a business process).
% - Trace sigma = <e1, ..., en>; |sigma| = number of events.
% - Event e = (c, a, d): case id c, activity a, data payload d (start/complete
%   timestamps, resource r, case attributes).
% - Projection functions act: E -> A and res: E -> R.
% - Sets A (activities) and R (resources) of the process.
% - Distinction event / activity / case; note that "event" and "activity
%   label" are often used interchangeably.
%
% subsec:bg:pm:prefix
% - k-prefix: sigma(k) = <e1, ..., ek> with k < |sigma|.
% - Running trace / ongoing execution: prefix of a trace that is not yet
%   complete; it is the input of the prescriptive problem.
% - Concept of "completion" of a prefix (its unknown suffix).
%
% subsec:bg:pm:ts
% - Why an abstraction is needed: avoid an over-deterministic transition
%   system and a state-space explosion.
% - Sequence-based state representation with window w:
%   S^w(sigma(k)) = last w activities of the prefix (or all of them if k < w).
% - Transition system mined from the log: maps abstract state -> activities
%   that can follow it.
% - Small example (reuse Fig. 2 of the paper or build one).


% #############################################################################
\section{Predictive Process Monitoring}
\label{sec:bg:ppm}
% #############################################################################
% Goal: how a property of the complete trace is predicted from a prefix;
% introduce the two targets used (outcome and remaining time) and the
% encoding.
% References: \cite{difrancescomarino2022} (PPM, Process Mining Handbook);
% \cite{marquezchamorro2018} (survey); \cite{teinemaa2018} (outcome-oriented);
% \cite{galanti2023}.
% Do NOT put here: your concrete sklearn pipeline, get_features() ->
% Methodology.
% Bridge: "PPM provides the estimate; Prescriptive Analytics uses it to
% decide what to do."

\subsection{The Predictive Process Monitoring Problem}
\label{subsec:bg:ppm:problem}
% - Definition: a function that, given a prefix, estimates a property of the
%   complete trace (outcome, remaining time, next activity, cost...).
% - Tasks relevant to the thesis:
%    * CASE OUTCOME: binary classification; defined by the occurrence of a key
%      activity (e.g. "O_Accepted" for BPIC2012/2017; "Network /
%      Back-Office Adjustment Requested" as a negative outcome for BAC).
%      The activity that defines the outcome is excluded from the
%      prescribable activities (fair evaluation).
%    * REMAINING TIME: regression on the time left until the case ends.
% - Notion of KPI function F(sigma, i): the KPI value after the first i events.

\subsection{Labeling}
\label{subsec:bg:ppm:labeling}
% - How the prefix training set is built: for every trace and every step k an
%   example (prefix, target) is generated.
% - Outcome label (boolean) and time label (continuous).

\subsection{Prefix Encoding}
\label{subsec:bg:ppm:encoding}
% - Why it is needed: tabular ML techniques require a fixed-length numeric
%   vector.
% - ACTIVITY-FREQUENCY ENCODING: a vector counting the occurrences of each
%   activity in the prefix. Other known options (last-state, aggregation,
%   index-based, embedding) to mention briefly for context.
% - Auxiliary attributes: prefix duration, current-event duration (end -
%   start, or fallback end - previous end), weekday of the last timestamp.
% - LABEL ENCODING of the candidate next activity and next resource
%   (NEXT_ACTIVITY, NEXT_RESOURCE): they are model inputs, not outputs.
% - Train/inference consistency: the same encoder is applied to the running
%   execution online.
% - One-hot vs label encoding: anticipate that the encoding choice affects
%   epistemic uncertainty (an unseen category becomes an all-zero row that
%   every sub-ensemble agrees on) -> picked up again in
%   \ref{sec:bg:uncertainty}.

\subsection{Normalization of the Time Target}
\label{subsec:bg:ppm:normalization}
% (BASIC STATISTICS.)
% - Problem: the remaining time has a strongly skewed (long-tail) distribution
%   and lives on a different scale than the binary outcome.
% - Three-step transformation (define it formally):
%    1. z-normalization (standardization): (x - mu) / s.
%    2. sigmoid: 1 / (1 + e^-z).
%    3. min-max scaling to [0, 1].
%   Result: the "sigmoid_mm" variable, nearly uniform on [0,1], which makes
%   the two KPIs comparable and makes even small improvements visible.
% - Invertibility of the transformation (needed to bring simulated times back
%   to seconds during evaluation): min-max^-1 -> logit -> de-standardization,
%   with edge clipping.
% - Recall: sample mean / standard deviation, the logistic function, monotone
%   transformations and order preservation.


% #############################################################################
\section{Prescriptive Process Analytics}
\label{sec:bg:prpa}
% #############################################################################
% Goal: introduce prescriptive process analytics (vs the other analytics
% types and vs predictive monitoring), give its formal problem statement, and
% explain what a KPI is. Three subsections: concept + examples; formal
% definition; KPIs.
% References: \cite{ibm_prescriptive} and \cite{wikipedia_prescriptive}
% (analytics taxonomy, technology behind prescriptive analytics -- consider
% replacing Wikipedia with a stronger source); \cite{kubrak2022} (Quo vadis?);
% \cite{le2025} (next step recommendation); \cite{le2025balancing} (the
% framework this thesis extends).
% Do NOT put here: the feasibility module internals, the search strategies
% (exhaustive vs genetic), the diverse top-k selection -> methodology chapter.
% Planning notes for that relocated material are kept as comments at the end
% of this section.
% Bridge: "When the KPIs are more than one and in conflict, 'optimal' must be
% redefined" -> Section \ref{sec:bg:moo}.

\subsection{From Prediction to Prescription}
\label{subsec:bg:prpa:pred2presc}
Data analytics is commonly described as a progression of four questions of
increasing sophistication and business value~\cite{ibm_prescriptive}:
\emph{descriptive} analytics asks \emph{what happened}, \emph{diagnostic}
analytics asks \emph{why it happened}, \emph{predictive} analytics asks
\emph{what is likely to happen next}, and \emph{prescriptive} analytics asks
\emph{what should be done next}. The four are complementary, each building on
the output of the previous one, but they differ in what they deliver: the
first three explain or forecast, whereas only the last one recommends a
course of action and makes the implications of each option
explicit~\cite{wikipedia_prescriptive}. Technically, prescriptive analytics
combines historical and contextual data with business rules and with
mathematical and computational models drawn from statistics, machine learning
and operations research, in order to reason about objectives, constraints,
uncertainty and trade-offs at
once~\cite{wikipedia_prescriptive,ibm_prescriptive}. It is applied, for
example, to recommend a treatment from a patient's characteristics, to set
prices and plan promotions from a demand forecast, to schedule maintenance
before a machine fails, or to score a transaction for fraud
risk~\cite{ibm_prescriptive}.

When this paradigm is applied to the executions of a business process
recorded in an event log, it is called \emph{Prescriptive Process Analytics}
(PrPA), or prescriptive process monitoring~\cite{kubrak2022}. PrPA supports
\emph{ongoing} process executions: instead of only forecasting how a running
case will end, it suggests interventions that steer it toward a desired
result, for instance preventing an unfavorable outcome or reducing the cost
or the duration of the case. The difference from predictive process
monitoring is one of intent. Predictive monitoring passively estimates a
property of the running case, such as its outcome, its remaining time or its
most likely next activity; prescriptive monitoring uses those estimates to
actively decide which action improves a target performance measure, taking
the dependencies between the different aspects of the process into account.

The kind of action prescribed varies across the literature. Early work
recommended the next activity to execute by exploiting the event data payload
and simple counterfactual reasoning~\cite{groger2014}. As prediction-based
methods spread, attention moved to flagging cases at risk of an undesired
outcome and to deciding \emph{whether} and \emph{when} to
intervene~\cite{teinemaa2018,fahrenkrogpetersen2022}. More recent approaches
recommend \emph{how} to proceed using causal inference~\cite{bozorgi2023},
reinforcement learning~\cite{branchi2022,metzger2020,shoush2022} or
interpretable, white-box predictors~\cite{padella2022}, and increasingly
prescribe not only the next activity but also the resource that should carry
it out --- the \emph{next step recommendation}~\cite{le2025} --- sometimes
accounting for how a resource's performance changes with
experience~\cite{padella2024}. Large language models have also started to be
used to make such recommendations easier to interpret~\cite{kubrak2024}.

What almost all of these approaches share is that they optimize a
\emph{single} performance measure at a time. The few that consider several
process aspects fold them into one fixed formula with a predetermined meaning
(e.g.\ monetary profit), which does not let a decision-maker examine
alternative trade-offs~\cite{le2025balancing}. Treating multiple, possibly
conflicting measures explicitly is the setting of this thesis; the tools for
it are introduced in Section~\ref{sec:bg:moo}.

\subsection{Formal Problem Definition}
\label{subsec:bg:prpa:formal}
A performance measure is formalized as a \emph{Key Performance Indicator}
(KPI) function $F : \mathcal{E}^{*} \times \mathbb{N} \to \mathbb{R}$ that,
given a trace $\sigma$ and a position $1 \le i \le |\sigma|$, returns the
value $F(\sigma, i)$ of that indicator for $\sigma$ after its first $i$
events~\cite{le2025balancing}. Evaluating a KPI at an intermediate position,
and not only on the completed trace, is what makes it usable while a case is
still running.

The prescription itself is a \emph{next step}: a pair
$(a, r) \in \mathcal{A} \times \mathcal{R}$ consisting of the activity to
perform next and the resource assigned to it. Given a running execution
$\sigma' = \langle e_1, \dots, e_k \rangle$, the prescriptive problem is to
choose the next event $e_{k+1}$ --- that is, the pair $(a, r)$ --- so that the
resulting completed trace
$\sigma_{\mathrm{pres}} = \langle e_1, \dots, e_k, e_{k+1}, \dots, e_n
\rangle$ attains the best possible KPI value. The candidate pairs are not
arbitrary: they are restricted to the activities that can actually follow the
current state of the case and to the resources able to perform them, a
feasibility constraint that is detailed, together with the rest of the
framework, in the methodology chapter. When a single KPI $F$ is considered,
``best'' simply means maximizing (or minimizing) $F$; when several KPIs
$F_1, \dots, F_n$ matter at once, they are aggregated into a single
\emph{utility function}
\begin{equation}
U_{\mathrm{KPI}}(\sigma, \boldsymbol{\lambda}, i)
  = \sum_{j=1}^{n} \lambda_j \, F_j(\sigma, i),
\label{eq:bg:prpa:utility}
\end{equation}
where $\boldsymbol{\lambda} = (\lambda_1, \dots, \lambda_n) \in [0, 1]^n$,
with $\sum_{j} \lambda_j = 1$, weights the relative importance of each
indicator according to the decision-maker's
priorities~\cite{le2025balancing}. The limitations of this linear
aggregation, and the alternative of not aggregating at all, are discussed in
Section~\ref{sec:bg:moo}.

\subsection{Key Performance Indicators}
\label{subsec:bg:prpa:kpi}
The objective of PrPA is stated informally as reaching a ``desired
outcome'', and a KPI is what turns that vague notion into a quantity that can
be optimized. Which quantity is appropriate depends entirely on the process
and on the organization's goals; recurring dimensions are the throughput time
of a case, its cost, some notion of quality or of regulatory compliance, and
the efficiency with which resources are used.

Two examples make the mapping concrete. In a loan-application process, a bank
that wants to grant more loans can take the \emph{acceptance rate} --- the
fraction of cases that reach a positive decision --- as a KPI to maximize,
whereas a bank that wants faster decisions can take the \emph{time to
decision} as a KPI to minimize. In a hospital, management may want to shorten
the \emph{length of stay} of patients in order to free capacity, while also
keeping the \emph{30-day readmission rate} low, so that discharging patients
earlier does not degrade the quality of care.

The two hospital KPIs illustrate a recurring difficulty: performance measures
often \emph{conflict}, so that improving one worsens another, and the right
balance between them is not fixed --- it depends on the domain, on the
regulatory context and on situational priorities. A prescription that is
optimal for one KPI can therefore be a poor choice overall. Making these
trade-offs explicit, instead of hiding them inside a single composite score,
is the problem addressed in Section~\ref{sec:bg:moo}.

% -----------------------------------------------------------------------------
% Planning notes -- material relocated to the methodology chapter, kept here
% for reference.
% -----------------------------------------------------------------------------
% Feasibility module:
% - Purpose: restrict the space to what is practically executable.
% - Activity filter: A|sigma(k) = activities that in the log follow the
%   abstract state of the prefix at least once (transition system, window w).
%   Fall back to progressively shorter suffixes if the exact prefix was never
%   seen.
% - Resource filter: R_a = resources that executed a at least once in the log.
% - Feasible prescription space: P_sigma(k) = {(a, r) : a in A|sigma(k),
%   r in R_a}.
% - Exclusion of "forbidden" activities (those defining the outcome).
%
% Search-space exploration strategies:
% - Exhaustive search: evaluates all feasible pairs; optimal w.r.t. the model
%   but costly when |P_sigma(k)| is large (many resources).
% - Heuristic / metaheuristic (genetic) methods: explore a subset.
%   (Algorithm detail belongs to Section \ref{sec:bg:moo} / methodology.)
%
% Extended single-KPI survey (if a separate Related Work chapter is wanted):
% - Resource allocation for global KPI improvement \cite{meneghello2024}.
% - de Leoni et al. on guidance for process continuation \cite{deleoni2020}.


% #############################################################################
\section{Conflicting KPIs and Multi-Objective Optimization}
\label{sec:bg:moo}
% #############################################################################
% Goal: the tools to understand why a Pareto front is used, how it is explored
% (enumeration, NSGA-II) and how a diversified subset of solutions is chosen.
% References: \cite{marler2004} (multi-objective survey); \cite{deb2002}
% (NSGA-II); \cite{mitchell1998} (genetic algorithms); \cite{mothilal2020}
% (DiCE); \cite{buliga2025}; p-dispersion / facility-location literature.
% Do NOT put here: the layout of the recommendations_*.csv files and the
% _objectives.csv sidecars -> Methodology / Evaluation.
% Bridge: "The quality of these decisions depends entirely on the quality -
% and the reliability - of the model that estimates the three objectives."

So far, most prescriptive-monitoring approaches optimize a single KPI. In
practice this is rarely enough: a business process is judged on several
performance measures at once, these measures are usually interdependent, and
improving one of them often comes at the expense of another. What is needed is
a way to reason about the whole set of measures together and to choose a
deliberate trade-off among them.

\subsection{Conflicting KPIs}
\label{subsec:bg:moo:devil}
The tension is easiest to see on an example. A hospital may want to shorten
the length of stay of its patients, in order to free capacity, and at the
same time keep the 30-day readmission rate low; discharging patients sooner
helps the first measure but tends to hurt the second, so no policy is best on
both counts and management has to settle on an acceptable compromise. The
same situation recurs across domains, and it is not limited to two measures:
the classical \emph{Devil's Quadrangle} already identifies four dimensions of
process performance --- time, cost, quality and flexibility --- that
typically cannot all be improved at once~\cite{dumas2018}.

\subsection{Scalarization and Its Limitations}
\label{subsec:bg:moo:scalarization}
One way to keep the convenience of a single optimum is to collapse the
several KPIs into one number. This is what the utility function of
Equation~\eqref{eq:bg:prpa:utility} does: a weighted sum of the individual
indicators, with weights $\boldsymbol{\lambda}$ chosen by the
decision-maker~\cite{marler2004}. Scalarization of this kind is simple and it
always yields one well-defined solution, but it has three drawbacks. It
forces the weights to be fixed \emph{before} any solution is seen, even
though the trade-off they encode is exactly what the decision-maker would
like to inspect; it returns a single point for each weight vector, hiding the
shape of the trade-off; and it requires the KPIs to be expressed on
comparable scales, which is why the time target is normalized as described in
Section~\ref{subsec:bg:ppm:normalization}.

\subsection{Pareto Optimality}
\label{subsec:bg:moo:pareto}
When the objectives genuinely conflict there is no single solution that is
best on all of them, and the problem instead has a whole set of equally
defensible compromises. This is captured by the notion of \emph{Pareto
optimality}~\cite{pareto1906,miettinen1999}. Consider a minimization problem
with $m$ objective functions $f_1, \dots, f_m$ and a decision vector $x$
ranging over a feasible set $X$. A solution $x_1$ is said to \emph{dominate}
a solution $x_2$, written $x_1 \prec x_2$, when it is no worse on every
objective and strictly better on at least one:
\begin{equation}
x_1 \prec x_2
\iff
\bigl( f_i(x_1) \le f_i(x_2)\ \ \forall\, i \in \{1, \dots, m\} \bigr)
\ \wedge\
\bigl( \exists\, j : f_j(x_1) < f_j(x_2) \bigr).
\label{eq:bg:moo:dominance}
\end{equation}
A solution is \emph{Pareto-optimal} --- equivalently \emph{efficient} or
\emph{non-dominated} --- if no feasible solution dominates it. The image of
the Pareto-optimal set in objective space is the \emph{Pareto front}
\begin{equation}
P^{*} = \bigl\{\, f(x) \;\big|\; x \in X,\ \nexists\, x' \in X :
f(x') \prec f(x) \,\bigr\}.
\label{eq:bg:moo:front}
\end{equation}
A point lies on the front simply because nothing dominates it; it need not
itself dominate any other point. The best and the worst value that each
objective attains over $P^{*}$ define the \emph{ideal} and the \emph{nadir}
point, which are convenient references for normalizing the objectives and for
measuring how close a given solution is to the (generally unattainable)
ideal. When two entire fronts have to be compared rather than two individual
points, a common scalar summary is the \emph{hypervolume}: the volume of
objective space that a front dominates with respect to a fixed reference
point~\cite{deb2001}.

Compared with scalarization, computing the Pareto front defers the trade-off
decision. The optimizer returns all the non-dominated compromises, and the
selection of one of them is made afterwards, with the cost of each
alternative in full view~\cite{honegumi}. In this thesis the front is built
on three objectives: the probability of a positive outcome, to be maximized;
the predicted remaining time, to be minimized; and the predictive standard
deviation of that time estimate, also to be minimized. The third objective
treats the model's own \emph{confidence} as a KPI, a choice motivated in
Section~\ref{sec:bg:uncertainty}.

\subsection{Front-Exploration Algorithms}
\label{subsec:bg:moo:algorithms}
% - Exact enumeration + front extraction: exact on small discrete spaces.
% - GENETIC ALGORITHMS: population, fitness, selection, crossover, mutation,
%   elitism, stopping criterion.
%    * The original paper's "custom fitness" approach: validity (val) +
%      alpha * proximity (prox), inspired by \cite{buliga2025}; DiCE as the
%      library \cite{mothilal2020}.
%    * NSGA-II \cite{deb2002}: non-dominated sorting into fronts, crowding
%      distance for diversity, tournament selection, elitism.
%      Implemented with pymoo (an integer variable indexing the (a,r) pair,
%      SBX + polynomial mutation + rounding repair), currently "parked"
%      because it is slower on small candidate sets: explain it anyway, it is
%      a methodological contribution.
% - Enumeration vs metaheuristic comparison: quality vs cost; robustness to
%   predictor errors on rare pairs (the paper's hypothesis: proximity favors
%   frequent pairs, which are estimated better).

\subsection{Selecting a Representative Subset of the Front}
\label{subsec:bg:moo:dispersion}
For a discrete problem the Pareto front can contain a large number of points,
some of them very close to one another. A recommender that has to present
alternatives to a user is often better served by a handful of clearly
distinct options than by the entire front. Choosing the $k$ most
representative points can be framed as a \emph{max-min dispersion}, or
\emph{$p$-dispersion}, problem: given $n$ candidate points with pairwise
distances $d_{ij}$, select a subset $S$ of size $k$ that maximizes the
smallest distance between any two selected points,
\begin{equation}
\max_{\substack{S \subseteq \{1, \dots, n\} \\ |S| = k}}\ \
\min_{\substack{i, j \in S \\ i \neq j}}\ d_{ij}.
\label{eq:bg:moo:dispersion}
\end{equation}
The problem was first studied in facility location, where the aim is to place
$p$ facilities so that the closest pair is as far apart as
possible~\cite{kuby1987}, and it has since been used more generally to
extract a maximally spread-out subset from a set of
candidates~\cite{maliszewski2012}. Because it is stated purely in terms of a
matrix of pairwise distances, the distance metric is a modeling choice; here
the natural one is the Euclidean distance between points in the normalized
objective space. The problem admits an exact mixed-integer programming
formulation, for which off-the-shelf solvers are
available~\cite{feng2022spopt}; the formulation actually used, together with
the rule that picks a single final recommendation, is given in the
methodology chapter.

Every objective on the front is produced by a predictive model, so the value
of the whole procedure rests on how accurate --- and how well calibrated ---
that model is. The next two sections describe it.


% #############################################################################
\section{Gradient Boosting and CatBoost}
\label{sec:bg:gbm}
% #############################################################################
% Goal: the model class used for the two predictors and what hyperparameters
% are / what role they play, plus their optimization. NO concrete values here.
% References: \cite{friedman2001} (gradient boosting); \cite{prokhorenkova2018}
% (CatBoost); \cite{akiba2019} (Optuna); \cite{bergstra2011} (TPE);
% \cite{galanti2023}.
% Do NOT put here: the search_spaces of 2_training_predictive_model.py, the
% number of trials, the best_hyperparams_*.json -> Experimental Setup (with a
% pointer to \ref{sec:bg:gbm}).
% Bridge: "A point model does not say how sure it is of what it predicts:
% uncertainty must be quantified."

\subsection{Tree Ensembles and Boosting}
\label{subsec:bg:gbm:boosting}
% - Decision / regression tree: splits, leaves, depth.
% - Bagging vs boosting (general idea).
% - Gradient boosting: sequential addition of trees fitting the negative
%   gradient of the loss; learning rate as shrinkage; number of trees
%   (iterations) and overfitting.

\subsection{CatBoost}
\label{subsec:bg:gbm:catboost}
% - Motivation for the choice: high predictive quality with contained
%   training time; native handling of categorical features; used in prior PPM
%   work \cite{galanti2023}.
% - Distinctive elements:
%    * ordered boosting / ordered target statistics (mitigates target leakage
%      / prediction shift);
%    * oblivious (symmetric) trees: same split at every level -> balanced
%      trees, fast inference, regularizing effect;
%    * growth alternatives (Depthwise, Lossguide) and min_data_in_leaf.
% - Losses used in the thesis (as CONCEPTS):
%    * Logloss for outcome classification;
%    * RMSE for regression; and the RMSEWithUncertainty variant, which
%      predicts mean AND variance (bridge to \ref{sec:bg:uncertainty}).
% - scale_pos_weight / class weights for imbalance (the positive-case
%   percentages differ a lot across logs: e.g. BAC ~82.5%, BPIC2012 ~47.9%).

\subsection{Hyperparameters: What They Are and What They Control}
\label{subsec:bg:gbm:hyperparams}
% - Definition: a parameter NOT learned from the data, fixed before training,
%   governing model capacity and the optimization process.
% - Grouped by ROLE (explain the role, not the value):
%    * Capacity / complexity: depth, grow_policy, min_data_in_leaf.
%    * Learning step: learning_rate, iterations.
%    * Regularization: l2_leaf_reg.
%    * Stochasticity / sub-sampling: bootstrap_type (Bayesian/Bernoulli/MVS),
%      bagging_temperature, subsample, colsample_bylevel.
%    * Imbalance: scale_pos_weight.
% - The bias-variance trade-off is governed by these groups.

\subsection{Hyperparameter Optimization}
\label{subsec:bg:gbm:hpo}
% - Why not by hand: large space, non-linear interactions.
% - Validation: hold-out, k-fold cross-validation, and - for processes - a
%   TEMPORAL split (avoid leakage from the future). Note the internal
%   "running" validation (the most recent training cases as validation).
% - Early stopping: stop adding trees when the validation metric stops
%   improving.
% - Search: grid / random search vs BAYESIAN OPTIMIZATION.
%    * Optuna: TPE (Tree-structured Parzen Estimator) as the surrogate;
%      define-by-run search space.
%    * Pruning (MedianPruner): interrupt unpromising trials.
% - Selection of the winning trial on ONE consistent metric (Logloss / RMSE)
%   and a final refit on 100% of the training data with the chosen tree count.


% #############################################################################
\section{Predictive Uncertainty Quantification}
\label{sec:bg:uncertainty}
% #############################################################################
% Goal: this is the THEORETICAL CORE of the contribution. Taxonomy of
% uncertainty, its estimation in ensembles and in gradient boosting,
% calibration. It must be the most developed section of the chapter.
% References: \cite{malinin2021} (Uncertainty in Gradient Boosting via
% Ensembles, ICLR 2021); \cite{hullermeier2021} (aleatoric vs epistemic,
% survey); \cite{lakshminarayanan2017} (deep ensembles); \cite{guo2017}
% (temperature scaling); \cite{kuleshov2018} (calibrated regression);
% \cite{angelopoulos2023} (conformal prediction); \cite{gal2016} (MC-dropout).
% Do NOT put here: the UncertaintyRegressor / UncertaintyClassifier classes,
% how the wrappers plug into the Pipeline, how out["std"] becomes objective #3
% -> Methodology. The obtained ECE/coverage/T/s values -> Evaluation.
% Bridge: "The Pareto front and the recommendations are evaluated with a
% process simulation, because a real A/B test is not possible."

\subsection{Taxonomy of Uncertainty}
\label{subsec:bg:uncertainty:taxonomy}
% - ALEATORIC uncertainty (data uncertainty): intrinsic / irreducible noise in
%   the data-generating process; does not shrink with more data.
%    * homoscedastic vs heteroscedastic.
% - EPISTEMIC uncertainty (knowledge / model uncertainty): the model's
%   ignorance due to lack of data / coverage; reducible with more data; high
%   on rare or out-of-distribution inputs.
% - TOTAL predictive uncertainty: the combination of the two.
% - Why distinguish them in the thesis: epistemic flags out-of-domain inputs /
%   rarely observed (a,r) pairs; aleatoric says how noisy the estimate is even
%   in the ideal case.
% - RESULT PREVIEW: on one-hot tabular data the regressor's epistemic part is
%   very small (~0.1% of the variance) and does NOT rise on never-seen inputs
%   (an unknown category -> an all-zero row every sub-ensemble agrees on). It
%   is reported anyway: the full decomposition is a result in itself
%   ("epistemic is negligible here, and why").

\subsection{Ensemble Methods for Uncertainty}
\label{subsec:bg:uncertainty:ensembles}
% - Idea: train/extract several models; the DISAGREEMENT among members
%   estimates the epistemic part, the AVERAGE member variance estimates the
%   aleatoric part.
% - Deep ensembles, MC-dropout, bootstrap ensembles: cite as context.
% - LAW OF TOTAL VARIANCE (regression):
%     Var[y] = E_m[ Var(y | m) ]  +  Var_m[ E(y | m) ]
%              \____ aleatoric ___/    \___ epistemic ___/
% - ENTROPY decomposition (classification):
%     H[ E_m p_m ]            = total uncertainty
%     E_m[ H(p_m) ]           = aleatoric uncertainty (expected entropy)
%     MI = total - aleatoric  = epistemic uncertainty (mutual information, BALD)
%   All in nats.

\subsection{Uncertainty in Gradient Boosting}
\label{subsec:bg:uncertainty:gbm}
% (\cite{malinin2021}.)
% - SGLB - Stochastic Gradient Langevin Boosting: noise injected into the
%   updates -> the model (approximately) samples from a posterior over the
%   parameters. In CatBoost: posterior_sampling=True.
% - VIRTUAL ENSEMBLES: from a single SGLB model several "sub-models" are
%   obtained using growing prefixes of the tree sequence -> an ensemble at the
%   cost of a single training run. Parameter: virtual_ensembles_count.
% - Regression with the RMSEWithUncertainty loss: each sub-model predicts
%   (mu_m, sigma_m^2); via the law of total variance:
%     data_std      = sqrt( mean_m sigma_m^2 )     (aleatoric)
%     knowledge_std = sqrt( var_m mu_m )           (epistemic)
%     total_std     = sqrt( data_var + knowledge_var )
% - Classification: prediction_type="TotalUncertainty" returns the expected
%   entropy (aleatoric) and the total entropy; the mutual information
%   (epistemic) is the difference.

\subsection{Calibration}
\label{subsec:bg:uncertainty:calibration}
% - What a CALIBRATED (probabilistic) model is: among the events predicted
%   with probability p, roughly a fraction p actually occur.
% - Diagnostics:
%    * reliability diagram;
%    * ECE - Expected Calibration Error: |confidence - accuracy| averaged over
%      probability bins;
%    * for regression: interval COVERAGE (fraction of |y - mu| within
%      +-1 sigma / +-2 sigma; target ~68% / ~95% for a Gaussian) and
%      SHARPNESS (typical / median interval width).
% - POST-HOC correction methods:
%    * TEMPERATURE SCALING \cite{guo2017}: a single scalar T,
%      p_cal = sigmoid( logit(p) / T ); T > 1 "softens" over-confident
%      probabilities. Monotone -> it does not reorder candidates by
%      probability, but it changes the magnitudes feeding the normalization
%      and the front distances.
%    * SIGMA RECALIBRATION for regression: a scalar factor s that rescales
%      every standard-deviation component so the empirical coverage matches
%      the nominal one; estimated on a held-out slice with a robust
%      median-of-standardized-residual estimator (avoids blow-ups when
%      CatBoost predicts a near-zero variance).
% - CONFORMAL PREDICTION as the RIGOROUS alternative \cite{angelopoulos2023}:
%   split-conformal on a dedicated holdout -> intervals with a distribution-
%   free marginal coverage guarantee. Cite it as "the rigorous alternative"
%   and motivate why the thesis used post-hoc recalibration instead (the
%   models are trained once, cost/benefit trade-off).

\subsection{Supporting Statistical Tools}
\label{subsec:bg:uncertainty:stats}
% (Compact recap.)
% - Variance and standard deviation as dispersion measures.
% - Shannon entropy of a Bernoulli; entropy as an uncertainty measure.
% - Mutual information / information gain.
% - The normal distribution, the 68-95-99.7 rule, the standardized residual,
%   the 0.6745 quantile of the normal (median absolute value).


% #############################################################################
\section{Business Process Simulation}
\label{sec:bg:simulation}
% #############################################################################
% Goal: what process simulation is, how a simulation model is obtained from
% the data, and with which metrics its fidelity is assessed.
% References: \cite{vinci2025prosit} (ProSiT); \cite{demoor2025} (SimBank);
% \cite{chapelacampa2024} (extraneous delays); \cite{vanderaalst2016} (Petri
% nets / soundness); \cite{vinci2025online} (online discovery).
% Do NOT put here: the prosit_simulation_results/... paths, the
% build_recommender_df code, the specific case ids -> Methodology / Evaluation.

\subsection{Why Simulation}
\label{subsec:bg:simulation:why}
% - Impossibility of an A/B test: one cannot intervene on a real process just
%   to evaluate.
% - COUNTERFACTUAL evaluation: simulate "what would have happened" following /
%   not following the recommendation.
% - Standard in PrPA for evaluating recommendations \cite{demoor2025}.

\subsection{Petri Nets and Process Models}
\label{subsec:bg:simulation:petri}
% - Petri net: places, transitions, arcs, marking, firing of a transition.
% - Workflow net and soundness (every case can always complete).
% - Initial / final marking; relation to the start/end of a case.
% - Replay and ALIGNMENT of a log prefix against the model: reconstruct the
%   marking reached by a historical prefix (log move, model move,
%   model-inserted activities).

\subsection{Stochastic Discrete-Event Simulation}
\label{subsec:bg:simulation:des}
% - Components of a process simulation model:
%    * case ARRIVAL process;
%    * transition branching weights / probabilities;
%    * activity EXECUTION times;
%    * WAITING times / queues;
%    * resource CALENDARS and availability;
%    * (extraneous delays \cite{chapelacampa2024}, as context).
% - Repeating the runs to mitigate stochasticity (e.g. 10 runs per scenario).

\subsection{Discovery of the Simulation Model}
\label{subsec:bg:simulation:discovery}
% - The simulation parameters are LEARNED from the historical log (resource
%   discovery, transition weights, time distributions, calendars).
% - METHODOLOGICAL NOTE: the simulator parameters are discovered from the FULL
%   log, not from the training split.
% - The Petri-net structure is a separate input (control-flow discovery),
%   distinct from the parameter discovery.
% - Online discovery for evolving processes \cite{vinci2025online} - context.

\subsection{ProSiT}
\label{subsec:bg:simulation:prosit}
% - The simulator used \cite{vinci2025prosit}: interactive and transparent
%   simulations.
% - What it returns: simulated logs; diagnostics (unreachable recommendations,
%   truncated "runaway" cases, non-fitting prefixes, model-inserted
%   activities).

\subsection{Assessing the Simulator's Fidelity}
\label{subsec:bg:simulation:fidelity}
% (STATISTICS.)
% - Real vs simulated distribution comparison:
%    * case-duration distributions (histograms / KDE, median);
%    * rate of positive-outcome cases, real vs simulated.
% - Distribution of the simulator's residuals.
% - (Optional) distribution-distance metrics: EMD / Wasserstein, KS; or a
%   comparison of summary statistics.

\subsection{KPI Improvement Metrics}
\label{subsec:bg:simulation:kpimetrics}
% (Definition; the numbers are in Evaluation.)
% - Delta_CO = average over prefixes of ( F_CO(sigma_pres) - F_CO(sigma_GT) ).
% - Delta_RT = average relative improvement of the remaining time w.r.t. the
%   ground truth.
% - Baseline / Ground Truth: no recommendation applied.
% - (Optional) confidence intervals on the deltas via bootstrap; 95%
%   confidence bands in the plots as lambda varies.


% #############################################################################
\section{Notation}
\label{sec:bg:notation}
% #############################################################################
% OPTIONAL but much appreciated by advisors. A half-page table of symbols,
% consistent with those used in the rest of the thesis.
% Uncomment and complete the table below (requires \usepackage{booktabs}).
%
% \begin{table}[h]
%   \centering
%   \begin{tabular}{ll}
%     \toprule
%     Symbol & Meaning \\
%     \midrule
%     $\mathcal{L}$              & event log \\
%     $\sigma,\ \sigma_{(k)}$    & trace, $k$-prefix \\
%     $\mathcal{A},\ \mathcal{R}$ & sets of activities and resources \\
%     $act,\ res$                & projection functions \\
%     $e = (c, a, \mathbf{d})$   & event (case id, activity, payload) \\
%     $S^{w}(\sigma_{(k)})$      & state abstraction with window $w$ \\
%     $\mathcal{P}_{\sigma_{(k)}}$ & feasible prescription space \\
%     $F_{CO},\ F_{RT}$          & KPIs: case outcome, remaining time (norm.) \\
%     $U_{KPI}$                  & utility function / combined KPI \\
%     $\lambda$                  & KPI trade-off weight \\
%     $\phi$                     & predictive model (regressor) \\
%     $p,\ p_{cal}$              & raw / calibrated outcome probability \\
%     $T$                        & calibration temperature \cite{guo2017} \\
%     $s$                        & std recalibration factor \\
%     $\sigma_{data}$            & aleatoric uncertainty (std) \\
%     $\sigma_{know}$            & epistemic uncertainty (std) \\
%     $\sigma_{tot}$             & total predictive uncertainty \\
%     $\mathrm{MI}$              & mutual information (epistemic, classif.) \\
%     $k$                        & n. of diversified recommendations per case \\
%     $\tau,\ \alpha$            & validity threshold, proximity weight \\
%     \bottomrule
%   \end{tabular}
%   \caption{Notation used in the thesis.}
%   \label{tab:notation}
% \end{table}


% =============================================================================
%  STATISTICAL TOOLS MAP -> WHERE THEY GO  (note to self, not for the thesis)
% =============================================================================
%   z-score, sigmoid, min-max, invertibility             -> \ref{subsec:bg:ppm:normalization}
%   Bayesian optimization / TPE, CV, temporal split,
%     early stopping, pruning                            -> \ref{subsec:bg:gbm:hpo}
%   variance / std, Shannon entropy, mutual information,
%     law of total variance                              -> \ref{subsec:bg:uncertainty:ensembles} / \ref{subsec:bg:uncertainty:stats}
%   ECE, reliability diagram, temperature scaling,
%     coverage, sharpness, conformal prediction          -> \ref{subsec:bg:uncertainty:calibration}
%   68-95-99.7 rule, standardized residual,
%     0.6745 quantile                                    -> \ref{subsec:bg:uncertainty:stats}
%   distribution comparison (KDE, median, EMD/KS),
%     residuals                                          -> \ref{subsec:bg:simulation:fidelity}
%   bootstrap / 95% confidence bands                     -> \ref{subsec:bg:simulation:kpimetrics} and Evaluation
%
%   HYPERPARAMETERS appear in Background ONLY as concepts
%   (\ref{subsec:bg:gbm:hyperparams}). Chosen values, search space, number of
%   Optuna trials and best-params tables -> Experimental Setup, with the
%   sentence "the concepts are introduced in Section~\ref{sec:bg:gbm}".
%
% =============================================================================
%  FINAL CHECKLIST BEFORE HANDING IN THE CHAPTER
% =============================================================================
%  [ ] Every section ends with a "bridge" to the next one.
%  [ ] No concrete experimental value in Background (concepts only).
%  [ ] No description of YOUR classes/files/functions (-> Methodology).
%  [ ] Every concept introduced is later used in Methodology or Evaluation
%      (no "orphan" background).
%  [ ] \ref{sec:bg:uncertainty} is the most developed section: it is the
%      contribution.
%  [ ] The aleatoric/epistemic taxonomy is explained BEFORE SGLB / virtual
%      ensembles.
%  [ ] Notation table consistent with the symbols in the rest of the thesis.
%  [ ] Every subsection has at least one solid bibliographic reference.
%  [ ] Agreed with the advisors whether Related Work stays here
%      (\ref{subsec:bg:prpa:sota}) or is a separate chapter.
% =============================================================================
 (or \include{...}).
%  - Packages expected from the main file: booktabs (notation table), amsmath,
%    hyperref, a bibliography backend (biblatex or bibtex).
%  - Replace the \cite{...} keys with the ones in your .bib file.
%
%  GUIDING PRINCIPLES (note to self, not for the thesis)
%  ----------------------------------------------------
%  - FUNNEL structure: from the general (process mining) to the specific
%    (uncertainty in gradient boosting). Each section introduces ONLY what the
%    next one needs.
%  - Follow the system PIPELINE:
%       offline (log -> encoding -> predictive model)
%    -> online (feasibility -> multi-objective search -> recommendation)
%    -> evaluation (business process simulation).
%  - Boundary with the other chapters:
%     * BACKGROUND  = WHAT a concept is and WHY it exists (reusable literature,
%                     independent of your work).
%     * METHODOLOGY = how YOU combine those concepts (architecture, the
%                     uncertainty wrappers, confidence as a Pareto objective).
%     * EVALUATION / EXPERIMENTAL SETUP = the concrete VALUES (Optuna search
%                     space, n. of trials, window_size, tau/alpha thresholds,
%                     ...) and the result tables.
%  - HYPERPARAMETERS: in Background you explain ONLY what a hyperparameter is
%    and what each group controls. The chosen values and the search space go
%    in the Experimental Setup, with a pointer to \ref{sec:bg:gbm}.
%  - STATISTICAL tools do NOT get a section of their own: spread them across
%    the point where the system uses them (see "STATISTICAL TOOLS MAP" at the
%    end of the file).
%
%  SECTION DEPENDENCY ORDER
%  -----------------------
%    PM --+--> PPM --+--> GBM/CatBoost --> Uncertainty
%         |          +--> PrPA --> Multi-objective opt.
%         +----------------------------> Simulation
%    If a topic is ONLY yours   -> Methodology.
%    If it is reusable literature -> Background (even if little known: NSGA-II,
%    SGLB, temperature scaling, p-dispersion, conformal prediction).
% =============================================================================


\chapter{Background}
\label{chap:background}

This chapter introduces the concepts on which the rest of the thesis builds.
The problem it addresses lies at the intersection of several fields: process
mining provides the data model (event logs and traces); predictive and
prescriptive process monitoring turn a running trace into a recommendation;
gradient boosting is the predictive engine, and quantifying its uncertainty
is the core of this work; multi-objective optimization frames the trade-off
between conflicting KPIs; and business process simulation is the vehicle used
to evaluate the recommendations.

% -----------------------------------------------------------------------------
% (Chapter introduction above. Original planning notes kept for reference.)
% -----------------------------------------------------------------------------
% Goal: tell the reader WHAT they need to know and WHY, previewing the
% "offline -> online -> evaluation" thread and flagging that the central
% section is the one on uncertainty (the contribution). Close with a sentence
% on the chapter's structure.
%
% Outline:
%  - The problem tackled by the thesis lives at the intersection of process
%    mining, machine learning and multi-objective optimization.
%  - Introduce the domain first (event log, PPM, PrPA), then the predictive
%    tool (gradient boosting) and its uncertainty, finally the process
%    simulation used for evaluation.
%  - List of the sections.


% #############################################################################
\section{Process Mining: Basic Concepts}
\label{sec:bg:pm}
% #############################################################################
% Goal of the section: fix the minimal vocabulary (log, trace, event, prefix,
% activity, resource) on which everything else rests.
% References: \cite{vanderaalst2016} (Process Mining: Data Science in Action);
% \cite{adams2021} (concept drift, to motivate the BPIC2017 split).
% Do NOT put here: your transition-system cache implementation
% (utils/setup_cache.py) -> Methodology.
% Bridge to the next section: "Predictive models are built on top of these
% prefixes."

Process mining is a discipline that combines data science and process
management: it analyzes the data recorded by information systems to
reconstruct how business processes actually run~\cite{vanderaalst2016}. The
basic idea is simple. Every time a process is executed, for example a
reimbursement request, an order, a patient's hospital journey, the systems
that support it leave traces in the form of event logs: they record which
activity was performed, when, and for which specific case.

From this data, process mining techniques make it possible to automatically
reconstruct the process model as it is actually executed, to compare it with
an expected model to identify deviations, and to enrich it with additional
information useful for improving it~\cite{vanderaalst2016}. In practice, it
shows organizations the difference between how they think their processes work
and how they actually work.

\subsection{Event Log, Traces, Events}
\label{subsec:bg:pm:log}
The starting point of any process-mining analysis is the \emph{event log}. An
event log $\mathcal{L}$ collects the executions of a business process, called
\emph{traces}: each trace $\sigma = \langle e_1, e_2, \dots, e_n \rangle$ is a
chronologically ordered sequence of events produced by one instance (or case)
of the process, and $|\sigma|$ denotes its length. Every event
$e = (c, a, \mathbf{d})$ is uniquely tied to a trace through a case identifier
$c$, records the execution of one activity $a$, and carries a vector of
additional attributes $\mathbf{d}$ such as the start
and completion timestamps of the activity and the resource that performed
it~\cite{le2025balancing}.
$\mathcal{A}$ and $\mathcal{R}$ denote, respectively, the set of activity
labels and the set of resources that occur anywhere in the process;
$\mathcal{E}$ is the universe of possible events and $\mathcal{E}^{*}$ the
space of finite event sequences over it. Two projection functions,
$act : \mathcal{E} \to \mathcal{A}$ and $res : \mathcal{E} \to \mathcal{R}$,
extract the activity and the resource of an event. Because an activity can
occur more than once along the same trace (e.g.\ inside a loop), ``event''
and ``activity occurrence'' are the more precise terms, even though the
shorter ``activity'' is used informally throughout whenever no ambiguity
arises.

\subsection{Prefixes and Running Traces}
\label{subsec:bg:pm:prefix}
Historical event logs only contain completed traces, but a prescriptive
system observes a process while it is still unfolding. Given a complete trace
$\sigma$, its $k$-\emph{prefix} $\sigma_{(k)} = \langle e_1, \dots, e_k
\rangle$, with $k < |\sigma|$, is the sequence of its first $k$
events~\cite{le2025balancing}; the
remaining $\langle e_{k+1}, \dots, e_n \rangle$ is the suffix that has not
happened yet. A \emph{running}, or ongoing, execution is exactly such a
prefix, produced by a case that has not reached completion.

\subsection{State Abstraction and Transition Systems}
\label{subsec:bg:pm:ts}
% (Needed by Section \ref{sec:bg:prpa}.)
A running execution, as introduced above, is itself a prefix. A common way
to reason about it in process mining is a \emph{transition
system}~\cite{vanderaalst2016}, a
directed graph whose nodes are \emph{states} and whose edges are labeled
with activities: an edge from state $s$ to state $s'$ labeled $a$ means
that executing $a$ from $s$ leads to $s'$. States abstract prefixes
through a state-representation function $S$, mapping every prefix
$\sigma_{(k)}$ to some state $S(\sigma_{(k)})$; two prefixes that map to
the same state are treated as equivalent, even if they differ elsewhere.

Given a log of completed traces $\mathcal{L}$, the transition system is
mined by recording, for every state that occurs in the log, the
activities empirically observed to follow it. Formally, the set of
feasible next activities for a prefix $\sigma_{(k)}$
is~\cite{le2025balancing}
\begin{equation}
\mathcal{A}|_{\sigma_{(k)}} = \{a \in \mathcal{A} \mid \exists\, \sigma'
\in \mathcal{L}, \exists\, l < |\sigma'| : S(\tilde\sigma_{(l)}) =
S(\sigma_{(k)}), a = act(\sigma'_{(l+1)})\}.
\label{eq:bg:pm:feasible}
\end{equation}
That is, $a$ is feasible after $\sigma_{(k)}$ if some historical trace
reached the same state and was then followed by $a$. This is what lets
the transition system say something about a prefix it has never
literally seen, as long as its state was reached, at some point, by
another prefix in the log.

The choice of $S$ has direct consequences on the resulting transition
system. The most direct choice, using the prefix itself as the state
($S(\sigma_{(k)}) = \sigma_{(k)}$), offers little help: a state would
then correspond to one specific history, so nothing learned from one
case would transfer to another, and the number of states would grow
with every distinct prefix in the log. Some abstraction of the prefix
into a state is therefore needed, coarse enough to generalize across
similar histories, yet informative enough to keep the resulting
recommendation meaningful.

% -----------------------------------------------------------------------------
% Original planning notes for this section, kept for reference.
% -----------------------------------------------------------------------------
% subsec:bg:pm:log
% - Definition of event log L as a collection of traces (execution sequences
%   of a business process).
% - Trace sigma = <e1, ..., en>; |sigma| = number of events.
% - Event e = (c, a, d): case id c, activity a, data payload d (start/complete
%   timestamps, resource r, case attributes).
% - Projection functions act: E -> A and res: E -> R.
% - Sets A (activities) and R (resources) of the process.
% - Distinction event / activity / case; note that "event" and "activity
%   label" are often used interchangeably.
%
% subsec:bg:pm:prefix
% - k-prefix: sigma(k) = <e1, ..., ek> with k < |sigma|.
% - Running trace / ongoing execution: prefix of a trace that is not yet
%   complete; it is the input of the prescriptive problem.
% - Concept of "completion" of a prefix (its unknown suffix).
%
% subsec:bg:pm:ts
% - Why an abstraction is needed: avoid an over-deterministic transition
%   system and a state-space explosion.
% - Sequence-based state representation with window w:
%   S^w(sigma(k)) = last w activities of the prefix (or all of them if k < w).
% - Transition system mined from the log: maps abstract state -> activities
%   that can follow it.
% - Small example (reuse Fig. 2 of the paper or build one).


% #############################################################################
\section{Predictive Process Monitoring}
\label{sec:bg:ppm}
% #############################################################################
% Goal: how a property of the complete trace is predicted from a prefix;
% introduce the two targets used (outcome and remaining time) and the
% encoding.
% References: \cite{difrancescomarino2022} (PPM, Process Mining Handbook);
% \cite{marquezchamorro2018} (survey); \cite{teinemaa2018} (outcome-oriented);
% \cite{galanti2023}.
% Do NOT put here: your concrete sklearn pipeline, get_features() ->
% Methodology.
% Bridge: "PPM provides the estimate; Prescriptive Analytics uses it to
% decide what to do."

\subsection{The Predictive Process Monitoring Problem}
\label{subsec:bg:ppm:problem}
% - Definition: a function that, given a prefix, estimates a property of the
%   complete trace (outcome, remaining time, next activity, cost...).
% - Tasks relevant to the thesis:
%    * CASE OUTCOME: binary classification; defined by the occurrence of a key
%      activity (e.g. "O_Accepted" for BPIC2012/2017; "Network /
%      Back-Office Adjustment Requested" as a negative outcome for BAC).
%      The activity that defines the outcome is excluded from the
%      prescribable activities (fair evaluation).
%    * REMAINING TIME: regression on the time left until the case ends.
% - Notion of KPI function F(sigma, i): the KPI value after the first i events.

\subsection{Labeling}
\label{subsec:bg:ppm:labeling}
% - How the prefix training set is built: for every trace and every step k an
%   example (prefix, target) is generated.
% - Outcome label (boolean) and time label (continuous).

\subsection{Prefix Encoding}
\label{subsec:bg:ppm:encoding}
% - Why it is needed: tabular ML techniques require a fixed-length numeric
%   vector.
% - ACTIVITY-FREQUENCY ENCODING: a vector counting the occurrences of each
%   activity in the prefix. Other known options (last-state, aggregation,
%   index-based, embedding) to mention briefly for context.
% - Auxiliary attributes: prefix duration, current-event duration (end -
%   start, or fallback end - previous end), weekday of the last timestamp.
% - LABEL ENCODING of the candidate next activity and next resource
%   (NEXT_ACTIVITY, NEXT_RESOURCE): they are model inputs, not outputs.
% - Train/inference consistency: the same encoder is applied to the running
%   execution online.
% - One-hot vs label encoding: anticipate that the encoding choice affects
%   epistemic uncertainty (an unseen category becomes an all-zero row that
%   every sub-ensemble agrees on) -> picked up again in
%   \ref{sec:bg:uncertainty}.

\subsection{Normalization of the Time Target}
\label{subsec:bg:ppm:normalization}
% (BASIC STATISTICS.)
% - Problem: the remaining time has a strongly skewed (long-tail) distribution
%   and lives on a different scale than the binary outcome.
% - Three-step transformation (define it formally):
%    1. z-normalization (standardization): (x - mu) / s.
%    2. sigmoid: 1 / (1 + e^-z).
%    3. min-max scaling to [0, 1].
%   Result: the "sigmoid_mm" variable, nearly uniform on [0,1], which makes
%   the two KPIs comparable and makes even small improvements visible.
% - Invertibility of the transformation (needed to bring simulated times back
%   to seconds during evaluation): min-max^-1 -> logit -> de-standardization,
%   with edge clipping.
% - Recall: sample mean / standard deviation, the logistic function, monotone
%   transformations and order preservation.


% #############################################################################
\section{Prescriptive Process Analytics}
\label{sec:bg:prpa}
% #############################################################################
% Goal: introduce prescriptive process analytics (vs the other analytics
% types and vs predictive monitoring), give its formal problem statement, and
% explain what a KPI is. Three subsections: concept + examples; formal
% definition; KPIs.
% References: \cite{ibm_prescriptive} and \cite{wikipedia_prescriptive}
% (analytics taxonomy, technology behind prescriptive analytics -- consider
% replacing Wikipedia with a stronger source); \cite{kubrak2022} (Quo vadis?);
% \cite{le2025} (next step recommendation); \cite{le2025balancing} (the
% framework this thesis extends).
% Do NOT put here: the feasibility module internals, the search strategies
% (exhaustive vs genetic), the diverse top-k selection -> methodology chapter.
% Planning notes for that relocated material are kept as comments at the end
% of this section.
% Bridge: "When the KPIs are more than one and in conflict, 'optimal' must be
% redefined" -> Section \ref{sec:bg:moo}.

\subsection{From Prediction to Prescription}
\label{subsec:bg:prpa:pred2presc}
Data analytics is commonly described as a progression of four questions of
increasing sophistication and business value~\cite{ibm_prescriptive}:
\emph{descriptive} analytics asks \emph{what happened}, \emph{diagnostic}
analytics asks \emph{why it happened}, \emph{predictive} analytics asks
\emph{what is likely to happen next}, and \emph{prescriptive} analytics asks
\emph{what should be done next}. The four are complementary, each building on
the output of the previous one, but they differ in what they deliver: the
first three explain or forecast, whereas only the last one recommends a
course of action and makes the implications of each option
explicit~\cite{wikipedia_prescriptive}. Technically, prescriptive analytics
combines historical and contextual data with business rules and with
mathematical and computational models drawn from statistics, machine learning
and operations research, in order to reason about objectives, constraints,
uncertainty and trade-offs at
once~\cite{wikipedia_prescriptive,ibm_prescriptive}. It is applied, for
example, to recommend a treatment from a patient's characteristics, to set
prices and plan promotions from a demand forecast, to schedule maintenance
before a machine fails, or to score a transaction for fraud
risk~\cite{ibm_prescriptive}.

When this paradigm is applied to the executions of a business process
recorded in an event log, it is called \emph{Prescriptive Process Analytics}
(PrPA), or prescriptive process monitoring~\cite{kubrak2022}. PrPA supports
\emph{ongoing} process executions: instead of only forecasting how a running
case will end, it suggests interventions that steer it toward a desired
result, for instance preventing an unfavorable outcome or reducing the cost
or the duration of the case. The difference from predictive process
monitoring is one of intent. Predictive monitoring passively estimates a
property of the running case, such as its outcome, its remaining time or its
most likely next activity; prescriptive monitoring uses those estimates to
actively decide which action improves a target performance measure, taking
the dependencies between the different aspects of the process into account.

The kind of action prescribed varies across the literature. Early work
recommended the next activity to execute by exploiting the event data payload
and simple counterfactual reasoning~\cite{groger2014}. As prediction-based
methods spread, attention moved to flagging cases at risk of an undesired
outcome and to deciding \emph{whether} and \emph{when} to
intervene~\cite{teinemaa2018,fahrenkrogpetersen2022}. More recent approaches
recommend \emph{how} to proceed using causal inference~\cite{bozorgi2023},
reinforcement learning~\cite{branchi2022,metzger2020,shoush2022} or
interpretable, white-box predictors~\cite{padella2022}, and increasingly
prescribe not only the next activity but also the resource that should carry
it out --- the \emph{next step recommendation}~\cite{le2025} --- sometimes
accounting for how a resource's performance changes with
experience~\cite{padella2024}. Large language models have also started to be
used to make such recommendations easier to interpret~\cite{kubrak2024}.

What almost all of these approaches share is that they optimize a
\emph{single} performance measure at a time. The few that consider several
process aspects fold them into one fixed formula with a predetermined meaning
(e.g.\ monetary profit), which does not let a decision-maker examine
alternative trade-offs~\cite{le2025balancing}. Treating multiple, possibly
conflicting measures explicitly is the setting of this thesis; the tools for
it are introduced in Section~\ref{sec:bg:moo}.

\subsection{Formal Problem Definition}
\label{subsec:bg:prpa:formal}
A performance measure is formalized as a \emph{Key Performance Indicator}
(KPI) function $F : \mathcal{E}^{*} \times \mathbb{N} \to \mathbb{R}$ that,
given a trace $\sigma$ and a position $1 \le i \le |\sigma|$, returns the
value $F(\sigma, i)$ of that indicator for $\sigma$ after its first $i$
events~\cite{le2025balancing}. Evaluating a KPI at an intermediate position,
and not only on the completed trace, is what makes it usable while a case is
still running.

The prescription itself is a \emph{next step}: a pair
$(a, r) \in \mathcal{A} \times \mathcal{R}$ consisting of the activity to
perform next and the resource assigned to it. Given a running execution
$\sigma' = \langle e_1, \dots, e_k \rangle$, the prescriptive problem is to
choose the next event $e_{k+1}$ --- that is, the pair $(a, r)$ --- so that the
resulting completed trace
$\sigma_{\mathrm{pres}} = \langle e_1, \dots, e_k, e_{k+1}, \dots, e_n
\rangle$ attains the best possible KPI value. The candidate pairs are not
arbitrary: they are restricted to the activities that can actually follow the
current state of the case and to the resources able to perform them, a
feasibility constraint that is detailed, together with the rest of the
framework, in the methodology chapter. When a single KPI $F$ is considered,
``best'' simply means maximizing (or minimizing) $F$; when several KPIs
$F_1, \dots, F_n$ matter at once, they are aggregated into a single
\emph{utility function}
\begin{equation}
U_{\mathrm{KPI}}(\sigma, \boldsymbol{\lambda}, i)
  = \sum_{j=1}^{n} \lambda_j \, F_j(\sigma, i),
\label{eq:bg:prpa:utility}
\end{equation}
where $\boldsymbol{\lambda} = (\lambda_1, \dots, \lambda_n) \in [0, 1]^n$,
with $\sum_{j} \lambda_j = 1$, weights the relative importance of each
indicator according to the decision-maker's
priorities~\cite{le2025balancing}. The limitations of this linear
aggregation, and the alternative of not aggregating at all, are discussed in
Section~\ref{sec:bg:moo}.

\subsection{Key Performance Indicators}
\label{subsec:bg:prpa:kpi}
The objective of PrPA is stated informally as reaching a ``desired
outcome'', and a KPI is what turns that vague notion into a quantity that can
be optimized. Which quantity is appropriate depends entirely on the process
and on the organization's goals; recurring dimensions are the throughput time
of a case, its cost, some notion of quality or of regulatory compliance, and
the efficiency with which resources are used.

Two examples make the mapping concrete. In a loan-application process, a bank
that wants to grant more loans can take the \emph{acceptance rate} --- the
fraction of cases that reach a positive decision --- as a KPI to maximize,
whereas a bank that wants faster decisions can take the \emph{time to
decision} as a KPI to minimize. In a hospital, management may want to shorten
the \emph{length of stay} of patients in order to free capacity, while also
keeping the \emph{30-day readmission rate} low, so that discharging patients
earlier does not degrade the quality of care.

The two hospital KPIs illustrate a recurring difficulty: performance measures
often \emph{conflict}, so that improving one worsens another, and the right
balance between them is not fixed --- it depends on the domain, on the
regulatory context and on situational priorities. A prescription that is
optimal for one KPI can therefore be a poor choice overall. Making these
trade-offs explicit, instead of hiding them inside a single composite score,
is the problem addressed in Section~\ref{sec:bg:moo}.

% -----------------------------------------------------------------------------
% Planning notes -- material relocated to the methodology chapter, kept here
% for reference.
% -----------------------------------------------------------------------------
% Feasibility module:
% - Purpose: restrict the space to what is practically executable.
% - Activity filter: A|sigma(k) = activities that in the log follow the
%   abstract state of the prefix at least once (transition system, window w).
%   Fall back to progressively shorter suffixes if the exact prefix was never
%   seen.
% - Resource filter: R_a = resources that executed a at least once in the log.
% - Feasible prescription space: P_sigma(k) = {(a, r) : a in A|sigma(k),
%   r in R_a}.
% - Exclusion of "forbidden" activities (those defining the outcome).
%
% Search-space exploration strategies:
% - Exhaustive search: evaluates all feasible pairs; optimal w.r.t. the model
%   but costly when |P_sigma(k)| is large (many resources).
% - Heuristic / metaheuristic (genetic) methods: explore a subset.
%   (Algorithm detail belongs to Section \ref{sec:bg:moo} / methodology.)
%
% Extended single-KPI survey (if a separate Related Work chapter is wanted):
% - Resource allocation for global KPI improvement \cite{meneghello2024}.
% - de Leoni et al. on guidance for process continuation \cite{deleoni2020}.


% #############################################################################
\section{Conflicting KPIs and Multi-Objective Optimization}
\label{sec:bg:moo}
% #############################################################################
% Goal: the tools to understand why a Pareto front is used, how it is explored
% (enumeration, NSGA-II) and how a diversified subset of solutions is chosen.
% References: \cite{marler2004} (multi-objective survey); \cite{deb2002}
% (NSGA-II); \cite{mitchell1998} (genetic algorithms); \cite{mothilal2020}
% (DiCE); \cite{buliga2025}; p-dispersion / facility-location literature.
% Do NOT put here: the layout of the recommendations_*.csv files and the
% _objectives.csv sidecars -> Methodology / Evaluation.
% Bridge: "The quality of these decisions depends entirely on the quality -
% and the reliability - of the model that estimates the three objectives."

So far, most prescriptive-monitoring approaches optimize a single KPI. In
practice this is rarely enough: a business process is judged on several
performance measures at once, these measures are usually interdependent, and
improving one of them often comes at the expense of another. What is needed is
a way to reason about the whole set of measures together and to choose a
deliberate trade-off among them.

\subsection{Conflicting KPIs}
\label{subsec:bg:moo:devil}
The tension is easiest to see on an example. A hospital may want to shorten
the length of stay of its patients, in order to free capacity, and at the
same time keep the 30-day readmission rate low; discharging patients sooner
helps the first measure but tends to hurt the second, so no policy is best on
both counts and management has to settle on an acceptable compromise. The
same situation recurs across domains, and it is not limited to two measures:
the classical \emph{Devil's Quadrangle} already identifies four dimensions of
process performance --- time, cost, quality and flexibility --- that
typically cannot all be improved at once~\cite{dumas2018}.

\subsection{Scalarization and Its Limitations}
\label{subsec:bg:moo:scalarization}
One way to keep the convenience of a single optimum is to collapse the
several KPIs into one number. This is what the utility function of
Equation~\eqref{eq:bg:prpa:utility} does: a weighted sum of the individual
indicators, with weights $\boldsymbol{\lambda}$ chosen by the
decision-maker~\cite{marler2004}. Scalarization of this kind is simple and it
always yields one well-defined solution, but it has three drawbacks. It
forces the weights to be fixed \emph{before} any solution is seen, even
though the trade-off they encode is exactly what the decision-maker would
like to inspect; it returns a single point for each weight vector, hiding the
shape of the trade-off; and it requires the KPIs to be expressed on
comparable scales, which is why the time target is normalized as described in
Section~\ref{subsec:bg:ppm:normalization}.

\subsection{Pareto Optimality}
\label{subsec:bg:moo:pareto}
When the objectives genuinely conflict there is no single solution that is
best on all of them, and the problem instead has a whole set of equally
defensible compromises. This is captured by the notion of \emph{Pareto
optimality}~\cite{pareto1906,miettinen1999}. Consider a minimization problem
with $m$ objective functions $f_1, \dots, f_m$ and a decision vector $x$
ranging over a feasible set $X$. A solution $x_1$ is said to \emph{dominate}
a solution $x_2$, written $x_1 \prec x_2$, when it is no worse on every
objective and strictly better on at least one:
\begin{equation}
x_1 \prec x_2
\iff
\bigl( f_i(x_1) \le f_i(x_2)\ \ \forall\, i \in \{1, \dots, m\} \bigr)
\ \wedge\
\bigl( \exists\, j : f_j(x_1) < f_j(x_2) \bigr).
\label{eq:bg:moo:dominance}
\end{equation}
A solution is \emph{Pareto-optimal} --- equivalently \emph{efficient} or
\emph{non-dominated} --- if no feasible solution dominates it. The image of
the Pareto-optimal set in objective space is the \emph{Pareto front}
\begin{equation}
P^{*} = \bigl\{\, f(x) \;\big|\; x \in X,\ \nexists\, x' \in X :
f(x') \prec f(x) \,\bigr\}.
\label{eq:bg:moo:front}
\end{equation}
A point lies on the front simply because nothing dominates it; it need not
itself dominate any other point. The best and the worst value that each
objective attains over $P^{*}$ define the \emph{ideal} and the \emph{nadir}
point, which are convenient references for normalizing the objectives and for
measuring how close a given solution is to the (generally unattainable)
ideal. When two entire fronts have to be compared rather than two individual
points, a common scalar summary is the \emph{hypervolume}: the volume of
objective space that a front dominates with respect to a fixed reference
point~\cite{deb2001}.

Compared with scalarization, computing the Pareto front defers the trade-off
decision. The optimizer returns all the non-dominated compromises, and the
selection of one of them is made afterwards, with the cost of each
alternative in full view~\cite{honegumi}. In this thesis the front is built
on three objectives: the probability of a positive outcome, to be maximized;
the predicted remaining time, to be minimized; and the predictive standard
deviation of that time estimate, also to be minimized. The third objective
treats the model's own \emph{confidence} as a KPI, a choice motivated in
Section~\ref{sec:bg:uncertainty}.

\subsection{Front-Exploration Algorithms}
\label{subsec:bg:moo:algorithms}
% - Exact enumeration + front extraction: exact on small discrete spaces.
% - GENETIC ALGORITHMS: population, fitness, selection, crossover, mutation,
%   elitism, stopping criterion.
%    * The original paper's "custom fitness" approach: validity (val) +
%      alpha * proximity (prox), inspired by \cite{buliga2025}; DiCE as the
%      library \cite{mothilal2020}.
%    * NSGA-II \cite{deb2002}: non-dominated sorting into fronts, crowding
%      distance for diversity, tournament selection, elitism.
%      Implemented with pymoo (an integer variable indexing the (a,r) pair,
%      SBX + polynomial mutation + rounding repair), currently "parked"
%      because it is slower on small candidate sets: explain it anyway, it is
%      a methodological contribution.
% - Enumeration vs metaheuristic comparison: quality vs cost; robustness to
%   predictor errors on rare pairs (the paper's hypothesis: proximity favors
%   frequent pairs, which are estimated better).

\subsection{Selecting a Representative Subset of the Front}
\label{subsec:bg:moo:dispersion}
For a discrete problem the Pareto front can contain a large number of points,
some of them very close to one another. A recommender that has to present
alternatives to a user is often better served by a handful of clearly
distinct options than by the entire front. Choosing the $k$ most
representative points can be framed as a \emph{max-min dispersion}, or
\emph{$p$-dispersion}, problem: given $n$ candidate points with pairwise
distances $d_{ij}$, select a subset $S$ of size $k$ that maximizes the
smallest distance between any two selected points,
\begin{equation}
\max_{\substack{S \subseteq \{1, \dots, n\} \\ |S| = k}}\ \
\min_{\substack{i, j \in S \\ i \neq j}}\ d_{ij}.
\label{eq:bg:moo:dispersion}
\end{equation}
The problem was first studied in facility location, where the aim is to place
$p$ facilities so that the closest pair is as far apart as
possible~\cite{kuby1987}, and it has since been used more generally to
extract a maximally spread-out subset from a set of
candidates~\cite{maliszewski2012}. Because it is stated purely in terms of a
matrix of pairwise distances, the distance metric is a modeling choice; here
the natural one is the Euclidean distance between points in the normalized
objective space. The problem admits an exact mixed-integer programming
formulation, for which off-the-shelf solvers are
available~\cite{feng2022spopt}; the formulation actually used, together with
the rule that picks a single final recommendation, is given in the
methodology chapter.

Every objective on the front is produced by a predictive model, so the value
of the whole procedure rests on how accurate --- and how well calibrated ---
that model is. The next two sections describe it.


% #############################################################################
\section{Gradient Boosting and CatBoost}
\label{sec:bg:gbm}
% #############################################################################
% Goal: the model class used for the two predictors and what hyperparameters
% are / what role they play, plus their optimization. NO concrete values here.
% References: \cite{friedman2001} (gradient boosting); \cite{prokhorenkova2018}
% (CatBoost); \cite{akiba2019} (Optuna); \cite{bergstra2011} (TPE);
% \cite{galanti2023}.
% Do NOT put here: the search_spaces of 2_training_predictive_model.py, the
% number of trials, the best_hyperparams_*.json -> Experimental Setup (with a
% pointer to \ref{sec:bg:gbm}).
% Bridge: "A point model does not say how sure it is of what it predicts:
% uncertainty must be quantified."

\subsection{Tree Ensembles and Boosting}
\label{subsec:bg:gbm:boosting}
% - Decision / regression tree: splits, leaves, depth.
% - Bagging vs boosting (general idea).
% - Gradient boosting: sequential addition of trees fitting the negative
%   gradient of the loss; learning rate as shrinkage; number of trees
%   (iterations) and overfitting.

\subsection{CatBoost}
\label{subsec:bg:gbm:catboost}
% - Motivation for the choice: high predictive quality with contained
%   training time; native handling of categorical features; used in prior PPM
%   work \cite{galanti2023}.
% - Distinctive elements:
%    * ordered boosting / ordered target statistics (mitigates target leakage
%      / prediction shift);
%    * oblivious (symmetric) trees: same split at every level -> balanced
%      trees, fast inference, regularizing effect;
%    * growth alternatives (Depthwise, Lossguide) and min_data_in_leaf.
% - Losses used in the thesis (as CONCEPTS):
%    * Logloss for outcome classification;
%    * RMSE for regression; and the RMSEWithUncertainty variant, which
%      predicts mean AND variance (bridge to \ref{sec:bg:uncertainty}).
% - scale_pos_weight / class weights for imbalance (the positive-case
%   percentages differ a lot across logs: e.g. BAC ~82.5%, BPIC2012 ~47.9%).

\subsection{Hyperparameters: What They Are and What They Control}
\label{subsec:bg:gbm:hyperparams}
% - Definition: a parameter NOT learned from the data, fixed before training,
%   governing model capacity and the optimization process.
% - Grouped by ROLE (explain the role, not the value):
%    * Capacity / complexity: depth, grow_policy, min_data_in_leaf.
%    * Learning step: learning_rate, iterations.
%    * Regularization: l2_leaf_reg.
%    * Stochasticity / sub-sampling: bootstrap_type (Bayesian/Bernoulli/MVS),
%      bagging_temperature, subsample, colsample_bylevel.
%    * Imbalance: scale_pos_weight.
% - The bias-variance trade-off is governed by these groups.

\subsection{Hyperparameter Optimization}
\label{subsec:bg:gbm:hpo}
% - Why not by hand: large space, non-linear interactions.
% - Validation: hold-out, k-fold cross-validation, and - for processes - a
%   TEMPORAL split (avoid leakage from the future). Note the internal
%   "running" validation (the most recent training cases as validation).
% - Early stopping: stop adding trees when the validation metric stops
%   improving.
% - Search: grid / random search vs BAYESIAN OPTIMIZATION.
%    * Optuna: TPE (Tree-structured Parzen Estimator) as the surrogate;
%      define-by-run search space.
%    * Pruning (MedianPruner): interrupt unpromising trials.
% - Selection of the winning trial on ONE consistent metric (Logloss / RMSE)
%   and a final refit on 100% of the training data with the chosen tree count.


% #############################################################################
\section{Predictive Uncertainty Quantification}
\label{sec:bg:uncertainty}
% #############################################################################
% Goal: this is the THEORETICAL CORE of the contribution. Taxonomy of
% uncertainty, its estimation in ensembles and in gradient boosting,
% calibration. It must be the most developed section of the chapter.
% References: \cite{malinin2021} (Uncertainty in Gradient Boosting via
% Ensembles, ICLR 2021); \cite{hullermeier2021} (aleatoric vs epistemic,
% survey); \cite{lakshminarayanan2017} (deep ensembles); \cite{guo2017}
% (temperature scaling); \cite{kuleshov2018} (calibrated regression);
% \cite{angelopoulos2023} (conformal prediction); \cite{gal2016} (MC-dropout).
% Do NOT put here: the UncertaintyRegressor / UncertaintyClassifier classes,
% how the wrappers plug into the Pipeline, how out["std"] becomes objective #3
% -> Methodology. The obtained ECE/coverage/T/s values -> Evaluation.
% Bridge: "The Pareto front and the recommendations are evaluated with a
% process simulation, because a real A/B test is not possible."

\subsection{Taxonomy of Uncertainty}
\label{subsec:bg:uncertainty:taxonomy}
% - ALEATORIC uncertainty (data uncertainty): intrinsic / irreducible noise in
%   the data-generating process; does not shrink with more data.
%    * homoscedastic vs heteroscedastic.
% - EPISTEMIC uncertainty (knowledge / model uncertainty): the model's
%   ignorance due to lack of data / coverage; reducible with more data; high
%   on rare or out-of-distribution inputs.
% - TOTAL predictive uncertainty: the combination of the two.
% - Why distinguish them in the thesis: epistemic flags out-of-domain inputs /
%   rarely observed (a,r) pairs; aleatoric says how noisy the estimate is even
%   in the ideal case.
% - RESULT PREVIEW: on one-hot tabular data the regressor's epistemic part is
%   very small (~0.1% of the variance) and does NOT rise on never-seen inputs
%   (an unknown category -> an all-zero row every sub-ensemble agrees on). It
%   is reported anyway: the full decomposition is a result in itself
%   ("epistemic is negligible here, and why").

\subsection{Ensemble Methods for Uncertainty}
\label{subsec:bg:uncertainty:ensembles}
% - Idea: train/extract several models; the DISAGREEMENT among members
%   estimates the epistemic part, the AVERAGE member variance estimates the
%   aleatoric part.
% - Deep ensembles, MC-dropout, bootstrap ensembles: cite as context.
% - LAW OF TOTAL VARIANCE (regression):
%     Var[y] = E_m[ Var(y | m) ]  +  Var_m[ E(y | m) ]
%              \____ aleatoric ___/    \___ epistemic ___/
% - ENTROPY decomposition (classification):
%     H[ E_m p_m ]            = total uncertainty
%     E_m[ H(p_m) ]           = aleatoric uncertainty (expected entropy)
%     MI = total - aleatoric  = epistemic uncertainty (mutual information, BALD)
%   All in nats.

\subsection{Uncertainty in Gradient Boosting}
\label{subsec:bg:uncertainty:gbm}
% (\cite{malinin2021}.)
% - SGLB - Stochastic Gradient Langevin Boosting: noise injected into the
%   updates -> the model (approximately) samples from a posterior over the
%   parameters. In CatBoost: posterior_sampling=True.
% - VIRTUAL ENSEMBLES: from a single SGLB model several "sub-models" are
%   obtained using growing prefixes of the tree sequence -> an ensemble at the
%   cost of a single training run. Parameter: virtual_ensembles_count.
% - Regression with the RMSEWithUncertainty loss: each sub-model predicts
%   (mu_m, sigma_m^2); via the law of total variance:
%     data_std      = sqrt( mean_m sigma_m^2 )     (aleatoric)
%     knowledge_std = sqrt( var_m mu_m )           (epistemic)
%     total_std     = sqrt( data_var + knowledge_var )
% - Classification: prediction_type="TotalUncertainty" returns the expected
%   entropy (aleatoric) and the total entropy; the mutual information
%   (epistemic) is the difference.

\subsection{Calibration}
\label{subsec:bg:uncertainty:calibration}
% - What a CALIBRATED (probabilistic) model is: among the events predicted
%   with probability p, roughly a fraction p actually occur.
% - Diagnostics:
%    * reliability diagram;
%    * ECE - Expected Calibration Error: |confidence - accuracy| averaged over
%      probability bins;
%    * for regression: interval COVERAGE (fraction of |y - mu| within
%      +-1 sigma / +-2 sigma; target ~68% / ~95% for a Gaussian) and
%      SHARPNESS (typical / median interval width).
% - POST-HOC correction methods:
%    * TEMPERATURE SCALING \cite{guo2017}: a single scalar T,
%      p_cal = sigmoid( logit(p) / T ); T > 1 "softens" over-confident
%      probabilities. Monotone -> it does not reorder candidates by
%      probability, but it changes the magnitudes feeding the normalization
%      and the front distances.
%    * SIGMA RECALIBRATION for regression: a scalar factor s that rescales
%      every standard-deviation component so the empirical coverage matches
%      the nominal one; estimated on a held-out slice with a robust
%      median-of-standardized-residual estimator (avoids blow-ups when
%      CatBoost predicts a near-zero variance).
% - CONFORMAL PREDICTION as the RIGOROUS alternative \cite{angelopoulos2023}:
%   split-conformal on a dedicated holdout -> intervals with a distribution-
%   free marginal coverage guarantee. Cite it as "the rigorous alternative"
%   and motivate why the thesis used post-hoc recalibration instead (the
%   models are trained once, cost/benefit trade-off).

\subsection{Supporting Statistical Tools}
\label{subsec:bg:uncertainty:stats}
% (Compact recap.)
% - Variance and standard deviation as dispersion measures.
% - Shannon entropy of a Bernoulli; entropy as an uncertainty measure.
% - Mutual information / information gain.
% - The normal distribution, the 68-95-99.7 rule, the standardized residual,
%   the 0.6745 quantile of the normal (median absolute value).


% #############################################################################
\section{Business Process Simulation}
\label{sec:bg:simulation}
% #############################################################################
% Goal: what process simulation is, how a simulation model is obtained from
% the data, and with which metrics its fidelity is assessed.
% References: \cite{vinci2025prosit} (ProSiT); \cite{demoor2025} (SimBank);
% \cite{chapelacampa2024} (extraneous delays); \cite{vanderaalst2016} (Petri
% nets / soundness); \cite{vinci2025online} (online discovery).
% Do NOT put here: the prosit_simulation_results/... paths, the
% build_recommender_df code, the specific case ids -> Methodology / Evaluation.

\subsection{Why Simulation}
\label{subsec:bg:simulation:why}
% - Impossibility of an A/B test: one cannot intervene on a real process just
%   to evaluate.
% - COUNTERFACTUAL evaluation: simulate "what would have happened" following /
%   not following the recommendation.
% - Standard in PrPA for evaluating recommendations \cite{demoor2025}.

\subsection{Petri Nets and Process Models}
\label{subsec:bg:simulation:petri}
% - Petri net: places, transitions, arcs, marking, firing of a transition.
% - Workflow net and soundness (every case can always complete).
% - Initial / final marking; relation to the start/end of a case.
% - Replay and ALIGNMENT of a log prefix against the model: reconstruct the
%   marking reached by a historical prefix (log move, model move,
%   model-inserted activities).

\subsection{Stochastic Discrete-Event Simulation}
\label{subsec:bg:simulation:des}
% - Components of a process simulation model:
%    * case ARRIVAL process;
%    * transition branching weights / probabilities;
%    * activity EXECUTION times;
%    * WAITING times / queues;
%    * resource CALENDARS and availability;
%    * (extraneous delays \cite{chapelacampa2024}, as context).
% - Repeating the runs to mitigate stochasticity (e.g. 10 runs per scenario).

\subsection{Discovery of the Simulation Model}
\label{subsec:bg:simulation:discovery}
% - The simulation parameters are LEARNED from the historical log (resource
%   discovery, transition weights, time distributions, calendars).
% - METHODOLOGICAL NOTE: the simulator parameters are discovered from the FULL
%   log, not from the training split.
% - The Petri-net structure is a separate input (control-flow discovery),
%   distinct from the parameter discovery.
% - Online discovery for evolving processes \cite{vinci2025online} - context.

\subsection{ProSiT}
\label{subsec:bg:simulation:prosit}
% - The simulator used \cite{vinci2025prosit}: interactive and transparent
%   simulations.
% - What it returns: simulated logs; diagnostics (unreachable recommendations,
%   truncated "runaway" cases, non-fitting prefixes, model-inserted
%   activities).

\subsection{Assessing the Simulator's Fidelity}
\label{subsec:bg:simulation:fidelity}
% (STATISTICS.)
% - Real vs simulated distribution comparison:
%    * case-duration distributions (histograms / KDE, median);
%    * rate of positive-outcome cases, real vs simulated.
% - Distribution of the simulator's residuals.
% - (Optional) distribution-distance metrics: EMD / Wasserstein, KS; or a
%   comparison of summary statistics.

\subsection{KPI Improvement Metrics}
\label{subsec:bg:simulation:kpimetrics}
% (Definition; the numbers are in Evaluation.)
% - Delta_CO = average over prefixes of ( F_CO(sigma_pres) - F_CO(sigma_GT) ).
% - Delta_RT = average relative improvement of the remaining time w.r.t. the
%   ground truth.
% - Baseline / Ground Truth: no recommendation applied.
% - (Optional) confidence intervals on the deltas via bootstrap; 95%
%   confidence bands in the plots as lambda varies.


% #############################################################################
\section{Notation}
\label{sec:bg:notation}
% #############################################################################
% OPTIONAL but much appreciated by advisors. A half-page table of symbols,
% consistent with those used in the rest of the thesis.
% Uncomment and complete the table below (requires \usepackage{booktabs}).
%
% \begin{table}[h]
%   \centering
%   \begin{tabular}{ll}
%     \toprule
%     Symbol & Meaning \\
%     \midrule
%     $\mathcal{L}$              & event log \\
%     $\sigma,\ \sigma_{(k)}$    & trace, $k$-prefix \\
%     $\mathcal{A},\ \mathcal{R}$ & sets of activities and resources \\
%     $act,\ res$                & projection functions \\
%     $e = (c, a, \mathbf{d})$   & event (case id, activity, payload) \\
%     $S^{w}(\sigma_{(k)})$      & state abstraction with window $w$ \\
%     $\mathcal{P}_{\sigma_{(k)}}$ & feasible prescription space \\
%     $F_{CO},\ F_{RT}$          & KPIs: case outcome, remaining time (norm.) \\
%     $U_{KPI}$                  & utility function / combined KPI \\
%     $\lambda$                  & KPI trade-off weight \\
%     $\phi$                     & predictive model (regressor) \\
%     $p,\ p_{cal}$              & raw / calibrated outcome probability \\
%     $T$                        & calibration temperature \cite{guo2017} \\
%     $s$                        & std recalibration factor \\
%     $\sigma_{data}$            & aleatoric uncertainty (std) \\
%     $\sigma_{know}$            & epistemic uncertainty (std) \\
%     $\sigma_{tot}$             & total predictive uncertainty \\
%     $\mathrm{MI}$              & mutual information (epistemic, classif.) \\
%     $k$                        & n. of diversified recommendations per case \\
%     $\tau,\ \alpha$            & validity threshold, proximity weight \\
%     \bottomrule
%   \end{tabular}
%   \caption{Notation used in the thesis.}
%   \label{tab:notation}
% \end{table}


% =============================================================================
%  STATISTICAL TOOLS MAP -> WHERE THEY GO  (note to self, not for the thesis)
% =============================================================================
%   z-score, sigmoid, min-max, invertibility             -> \ref{subsec:bg:ppm:normalization}
%   Bayesian optimization / TPE, CV, temporal split,
%     early stopping, pruning                            -> \ref{subsec:bg:gbm:hpo}
%   variance / std, Shannon entropy, mutual information,
%     law of total variance                              -> \ref{subsec:bg:uncertainty:ensembles} / \ref{subsec:bg:uncertainty:stats}
%   ECE, reliability diagram, temperature scaling,
%     coverage, sharpness, conformal prediction          -> \ref{subsec:bg:uncertainty:calibration}
%   68-95-99.7 rule, standardized residual,
%     0.6745 quantile                                    -> \ref{subsec:bg:uncertainty:stats}
%   distribution comparison (KDE, median, EMD/KS),
%     residuals                                          -> \ref{subsec:bg:simulation:fidelity}
%   bootstrap / 95% confidence bands                     -> \ref{subsec:bg:simulation:kpimetrics} and Evaluation
%
%   HYPERPARAMETERS appear in Background ONLY as concepts
%   (\ref{subsec:bg:gbm:hyperparams}). Chosen values, search space, number of
%   Optuna trials and best-params tables -> Experimental Setup, with the
%   sentence "the concepts are introduced in Section~\ref{sec:bg:gbm}".
%
% =============================================================================
%  FINAL CHECKLIST BEFORE HANDING IN THE CHAPTER
% =============================================================================
%  [ ] Every section ends with a "bridge" to the next one.
%  [ ] No concrete experimental value in Background (concepts only).
%  [ ] No description of YOUR classes/files/functions (-> Methodology).
%  [ ] Every concept introduced is later used in Methodology or Evaluation
%      (no "orphan" background).
%  [ ] \ref{sec:bg:uncertainty} is the most developed section: it is the
%      contribution.
%  [ ] The aleatoric/epistemic taxonomy is explained BEFORE SGLB / virtual
%      ensembles.
%  [ ] Notation table consistent with the symbols in the rest of the thesis.
%  [ ] Every subsection has at least one solid bibliographic reference.
%  [ ] Agreed with the advisors whether Related Work stays here
%      (\ref{subsec:bg:prpa:sota}) or is a separate chapter.
% =============================================================================
 (or \include{...}).
%  - Packages expected from the main file: booktabs (notation table), amsmath,
%    hyperref, a bibliography backend (biblatex or bibtex).
%  - Replace the \cite{...} keys with the ones in your .bib file.
%
%  GUIDING PRINCIPLES (note to self, not for the thesis)
%  ----------------------------------------------------
%  - FUNNEL structure: from the general (process mining) to the specific
%    (uncertainty in gradient boosting). Each section introduces ONLY what the
%    next one needs.
%  - Follow the system PIPELINE:
%       offline (log -> encoding -> predictive model)
%    -> online (feasibility -> multi-objective search -> recommendation)
%    -> evaluation (business process simulation).
%  - Boundary with the other chapters:
%     * BACKGROUND  = WHAT a concept is and WHY it exists (reusable literature,
%                     independent of your work).
%     * METHODOLOGY = how YOU combine those concepts (architecture, the
%                     uncertainty wrappers, confidence as a Pareto objective).
%     * EVALUATION / EXPERIMENTAL SETUP = the concrete VALUES (Optuna search
%                     space, n. of trials, window_size, tau/alpha thresholds,
%                     ...) and the result tables.
%  - HYPERPARAMETERS: in Background you explain ONLY what a hyperparameter is
%    and what each group controls. The chosen values and the search space go
%    in the Experimental Setup, with a pointer to \ref{sec:bg:gbm}.
%  - STATISTICAL tools do NOT get a section of their own: spread them across
%    the point where the system uses them (see "STATISTICAL TOOLS MAP" at the
%    end of the file).
%
%  SECTION DEPENDENCY ORDER
%  -----------------------
%    PM --+--> PPM --+--> GBM/CatBoost --> Uncertainty
%         |          +--> PrPA --> Multi-objective opt.
%         +----------------------------> Simulation
%    If a topic is ONLY yours   -> Methodology.
%    If it is reusable literature -> Background (even if little known: NSGA-II,
%    SGLB, temperature scaling, p-dispersion, conformal prediction).
% =============================================================================


\chapter{Background}
\label{chap:background}

This chapter introduces the concepts on which the rest of the thesis builds.
The problem it addresses lies at the intersection of several fields: process
mining provides the data model (event logs and traces); predictive and
prescriptive process monitoring turn a running trace into a recommendation;
gradient boosting is the predictive engine, and quantifying its uncertainty
is the core of this work; multi-objective optimization frames the trade-off
between conflicting KPIs; and business process simulation is the vehicle used
to evaluate the recommendations.

% -----------------------------------------------------------------------------
% (Chapter introduction above. Original planning notes kept for reference.)
% -----------------------------------------------------------------------------
% Goal: tell the reader WHAT they need to know and WHY, previewing the
% "offline -> online -> evaluation" thread and flagging that the central
% section is the one on uncertainty (the contribution). Close with a sentence
% on the chapter's structure.
%
% Outline:
%  - The problem tackled by the thesis lives at the intersection of process
%    mining, machine learning and multi-objective optimization.
%  - Introduce the domain first (event log, PPM, PrPA), then the predictive
%    tool (gradient boosting) and its uncertainty, finally the process
%    simulation used for evaluation.
%  - List of the sections.


% #############################################################################
\section{Process Mining: Basic Concepts}
\label{sec:bg:pm}
% #############################################################################
% Goal of the section: fix the minimal vocabulary (log, trace, event, prefix,
% activity, resource) on which everything else rests.
% References: \cite{vanderaalst2016} (Process Mining: Data Science in Action);
% \cite{adams2021} (concept drift, to motivate the BPIC2017 split).
% Do NOT put here: your transition-system cache implementation
% (utils/setup_cache.py) -> Methodology.
% Bridge to the next section: "Predictive models are built on top of these
% prefixes."

Process mining is a discipline that combines data science and process
management: it analyzes the data recorded by information systems to
reconstruct how business processes actually run~\cite{vanderaalst2016}. The
basic idea is simple. Every time a process is executed, for example a
reimbursement request, an order, a patient's hospital journey, the systems
that support it leave traces in the form of event logs: they record which
activity was performed, when, and for which specific case.

From this data, process mining techniques make it possible to automatically
reconstruct the process model as it is actually executed, to compare it with
an expected model to identify deviations, and to enrich it with additional
information useful for improving it~\cite{vanderaalst2016}. In practice, it
shows organizations the difference between how they think their processes work
and how they actually work.

\subsection{Event Log, Traces, Events}
\label{subsec:bg:pm:log}
The starting point of any process-mining analysis is the \emph{event log}. An
event log $\mathcal{L}$ collects the executions of a business process, called
\emph{traces}: each trace $\sigma = \langle e_1, e_2, \dots, e_n \rangle$ is a
chronologically ordered sequence of events produced by one instance (or case)
of the process, and $|\sigma|$ denotes its length. Every event
$e = (c, a, \mathbf{d})$ is uniquely tied to a trace through a case identifier
$c$, records the execution of one activity $a$, and carries a vector of
additional attributes $\mathbf{d}$ such as the start
and completion timestamps of the activity and the resource that performed
it~\cite{le2025balancing}.
$\mathcal{A}$ and $\mathcal{R}$ denote, respectively, the set of activity
labels and the set of resources that occur anywhere in the process;
$\mathcal{E}$ is the universe of possible events and $\mathcal{E}^{*}$ the
space of finite event sequences over it. Two projection functions,
$act : \mathcal{E} \to \mathcal{A}$ and $res : \mathcal{E} \to \mathcal{R}$,
extract the activity and the resource of an event. Because an activity can
occur more than once along the same trace (e.g.\ inside a loop), ``event''
and ``activity occurrence'' are the more precise terms, even though the
shorter ``activity'' is used informally throughout whenever no ambiguity
arises.

\subsection{Prefixes and Running Traces}
\label{subsec:bg:pm:prefix}
Historical event logs only contain completed traces, but a prescriptive
system observes a process while it is still unfolding. Given a complete trace
$\sigma$, its $k$-\emph{prefix} $\sigma_{(k)} = \langle e_1, \dots, e_k
\rangle$, with $k < |\sigma|$, is the sequence of its first $k$
events~\cite{le2025balancing}; the
remaining $\langle e_{k+1}, \dots, e_n \rangle$ is the suffix that has not
happened yet. A \emph{running}, or ongoing, execution is exactly such a
prefix, produced by a case that has not reached completion.

\subsection{State Abstraction and Transition Systems}
\label{subsec:bg:pm:ts}
% (Needed by Section \ref{sec:bg:prpa}.)
A running execution, as introduced above, is itself a prefix. A common way
to reason about it in process mining is a \emph{transition
system}~\cite{vanderaalst2016}, a
directed graph whose nodes are \emph{states} and whose edges are labeled
with activities: an edge from state $s$ to state $s'$ labeled $a$ means
that executing $a$ from $s$ leads to $s'$. States abstract prefixes
through a state-representation function $S$, mapping every prefix
$\sigma_{(k)}$ to some state $S(\sigma_{(k)})$; two prefixes that map to
the same state are treated as equivalent, even if they differ elsewhere.

Given a log of completed traces $\mathcal{L}$, the transition system is
mined by recording, for every state that occurs in the log, the
activities empirically observed to follow it. Formally, the set of
feasible next activities for a prefix $\sigma_{(k)}$
is~\cite{le2025balancing}
\begin{equation}
\mathcal{A}|_{\sigma_{(k)}} = \{a \in \mathcal{A} \mid \exists\, \sigma'
\in \mathcal{L}, \exists\, l < |\sigma'| : S(\tilde\sigma_{(l)}) =
S(\sigma_{(k)}), a = act(\sigma'_{(l+1)})\}.
\label{eq:bg:pm:feasible}
\end{equation}
That is, $a$ is feasible after $\sigma_{(k)}$ if some historical trace
reached the same state and was then followed by $a$. This is what lets
the transition system say something about a prefix it has never
literally seen, as long as its state was reached, at some point, by
another prefix in the log.

The choice of $S$ has direct consequences on the resulting transition
system. The most direct choice, using the prefix itself as the state
($S(\sigma_{(k)}) = \sigma_{(k)}$), offers little help: a state would
then correspond to one specific history, so nothing learned from one
case would transfer to another, and the number of states would grow
with every distinct prefix in the log. Some abstraction of the prefix
into a state is therefore needed, coarse enough to generalize across
similar histories, yet informative enough to keep the resulting
recommendation meaningful.

% -----------------------------------------------------------------------------
% Original planning notes for this section, kept for reference.
% -----------------------------------------------------------------------------
% subsec:bg:pm:log
% - Definition of event log L as a collection of traces (execution sequences
%   of a business process).
% - Trace sigma = <e1, ..., en>; |sigma| = number of events.
% - Event e = (c, a, d): case id c, activity a, data payload d (start/complete
%   timestamps, resource r, case attributes).
% - Projection functions act: E -> A and res: E -> R.
% - Sets A (activities) and R (resources) of the process.
% - Distinction event / activity / case; note that "event" and "activity
%   label" are often used interchangeably.
%
% subsec:bg:pm:prefix
% - k-prefix: sigma(k) = <e1, ..., ek> with k < |sigma|.
% - Running trace / ongoing execution: prefix of a trace that is not yet
%   complete; it is the input of the prescriptive problem.
% - Concept of "completion" of a prefix (its unknown suffix).
%
% subsec:bg:pm:ts
% - Why an abstraction is needed: avoid an over-deterministic transition
%   system and a state-space explosion.
% - Sequence-based state representation with window w:
%   S^w(sigma(k)) = last w activities of the prefix (or all of them if k < w).
% - Transition system mined from the log: maps abstract state -> activities
%   that can follow it.
% - Small example (reuse Fig. 2 of the paper or build one).


% #############################################################################
\section{Predictive Process Monitoring}
\label{sec:bg:ppm}
% #############################################################################
% Goal: how a property of the complete trace is predicted from a prefix;
% introduce the two targets used (outcome and remaining time) and the
% encoding.
% References: \cite{difrancescomarino2022} (PPM, Process Mining Handbook);
% \cite{marquezchamorro2018} (survey); \cite{teinemaa2018} (outcome-oriented);
% \cite{galanti2023}.
% Do NOT put here: your concrete sklearn pipeline, get_features() ->
% Methodology.
% Bridge: "PPM provides the estimate; Prescriptive Analytics uses it to
% decide what to do."

\subsection{The Predictive Process Monitoring Problem}
\label{subsec:bg:ppm:problem}
% - Definition: a function that, given a prefix, estimates a property of the
%   complete trace (outcome, remaining time, next activity, cost...).
% - Tasks relevant to the thesis:
%    * CASE OUTCOME: binary classification; defined by the occurrence of a key
%      activity (e.g. "O_Accepted" for BPIC2012/2017; "Network /
%      Back-Office Adjustment Requested" as a negative outcome for BAC).
%      The activity that defines the outcome is excluded from the
%      prescribable activities (fair evaluation).
%    * REMAINING TIME: regression on the time left until the case ends.
% - Notion of KPI function F(sigma, i): the KPI value after the first i events.

\subsection{Labeling}
\label{subsec:bg:ppm:labeling}
% - How the prefix training set is built: for every trace and every step k an
%   example (prefix, target) is generated.
% - Outcome label (boolean) and time label (continuous).

\subsection{Prefix Encoding}
\label{subsec:bg:ppm:encoding}
% - Why it is needed: tabular ML techniques require a fixed-length numeric
%   vector.
% - ACTIVITY-FREQUENCY ENCODING: a vector counting the occurrences of each
%   activity in the prefix. Other known options (last-state, aggregation,
%   index-based, embedding) to mention briefly for context.
% - Auxiliary attributes: prefix duration, current-event duration (end -
%   start, or fallback end - previous end), weekday of the last timestamp.
% - LABEL ENCODING of the candidate next activity and next resource
%   (NEXT_ACTIVITY, NEXT_RESOURCE): they are model inputs, not outputs.
% - Train/inference consistency: the same encoder is applied to the running
%   execution online.
% - One-hot vs label encoding: anticipate that the encoding choice affects
%   epistemic uncertainty (an unseen category becomes an all-zero row that
%   every sub-ensemble agrees on) -> picked up again in
%   \ref{sec:bg:uncertainty}.

\subsection{Normalization of the Time Target}
\label{subsec:bg:ppm:normalization}
% (BASIC STATISTICS.)
% - Problem: the remaining time has a strongly skewed (long-tail) distribution
%   and lives on a different scale than the binary outcome.
% - Three-step transformation (define it formally):
%    1. z-normalization (standardization): (x - mu) / s.
%    2. sigmoid: 1 / (1 + e^-z).
%    3. min-max scaling to [0, 1].
%   Result: the "sigmoid_mm" variable, nearly uniform on [0,1], which makes
%   the two KPIs comparable and makes even small improvements visible.
% - Invertibility of the transformation (needed to bring simulated times back
%   to seconds during evaluation): min-max^-1 -> logit -> de-standardization,
%   with edge clipping.
% - Recall: sample mean / standard deviation, the logistic function, monotone
%   transformations and order preservation.


% #############################################################################
\section{Prescriptive Process Analytics}
\label{sec:bg:prpa}
% #############################################################################
% Goal: introduce prescriptive process analytics (vs the other analytics
% types and vs predictive monitoring), give its formal problem statement, and
% explain what a KPI is. Three subsections: concept + examples; formal
% definition; KPIs.
% References: \cite{ibm_prescriptive} and \cite{wikipedia_prescriptive}
% (analytics taxonomy, technology behind prescriptive analytics -- consider
% replacing Wikipedia with a stronger source); \cite{kubrak2022} (Quo vadis?);
% \cite{le2025} (next step recommendation); \cite{le2025balancing} (the
% framework this thesis extends).
% Do NOT put here: the feasibility module internals, the search strategies
% (exhaustive vs genetic), the diverse top-k selection -> methodology chapter.
% Planning notes for that relocated material are kept as comments at the end
% of this section.
% Bridge: "When the KPIs are more than one and in conflict, 'optimal' must be
% redefined" -> Section \ref{sec:bg:moo}.

\subsection{From Prediction to Prescription}
\label{subsec:bg:prpa:pred2presc}
Data analytics is commonly described as a progression of four questions of
increasing sophistication and business value~\cite{ibm_prescriptive}:
\emph{descriptive} analytics asks \emph{what happened}, \emph{diagnostic}
analytics asks \emph{why it happened}, \emph{predictive} analytics asks
\emph{what is likely to happen next}, and \emph{prescriptive} analytics asks
\emph{what should be done next}. The four are complementary, each building on
the output of the previous one, but they differ in what they deliver: the
first three explain or forecast, whereas only the last one recommends a
course of action and makes the implications of each option
explicit~\cite{wikipedia_prescriptive}. Technically, prescriptive analytics
combines historical and contextual data with business rules and with
mathematical and computational models drawn from statistics, machine learning
and operations research, in order to reason about objectives, constraints,
uncertainty and trade-offs at
once~\cite{wikipedia_prescriptive,ibm_prescriptive}. It is applied, for
example, to recommend a treatment from a patient's characteristics, to set
prices and plan promotions from a demand forecast, to schedule maintenance
before a machine fails, or to score a transaction for fraud
risk~\cite{ibm_prescriptive}.

When this paradigm is applied to the executions of a business process
recorded in an event log, it is called \emph{Prescriptive Process Analytics}
(PrPA), or prescriptive process monitoring~\cite{kubrak2022}. PrPA supports
\emph{ongoing} process executions: instead of only forecasting how a running
case will end, it suggests interventions that steer it toward a desired
result, for instance preventing an unfavorable outcome or reducing the cost
or the duration of the case. The difference from predictive process
monitoring is one of intent. Predictive monitoring passively estimates a
property of the running case, such as its outcome, its remaining time or its
most likely next activity; prescriptive monitoring uses those estimates to
actively decide which action improves a target performance measure, taking
the dependencies between the different aspects of the process into account.

The kind of action prescribed varies across the literature. Early work
recommended the next activity to execute by exploiting the event data payload
and simple counterfactual reasoning~\cite{groger2014}. As prediction-based
methods spread, attention moved to flagging cases at risk of an undesired
outcome and to deciding \emph{whether} and \emph{when} to
intervene~\cite{teinemaa2018,fahrenkrogpetersen2022}. More recent approaches
recommend \emph{how} to proceed using causal inference~\cite{bozorgi2023},
reinforcement learning~\cite{branchi2022,metzger2020,shoush2022} or
interpretable, white-box predictors~\cite{padella2022}, and increasingly
prescribe not only the next activity but also the resource that should carry
it out --- the \emph{next step recommendation}~\cite{le2025} --- sometimes
accounting for how a resource's performance changes with
experience~\cite{padella2024}. Large language models have also started to be
used to make such recommendations easier to interpret~\cite{kubrak2024}.

What almost all of these approaches share is that they optimize a
\emph{single} performance measure at a time. The few that consider several
process aspects fold them into one fixed formula with a predetermined meaning
(e.g.\ monetary profit), which does not let a decision-maker examine
alternative trade-offs~\cite{le2025balancing}. Treating multiple, possibly
conflicting measures explicitly is the setting of this thesis; the tools for
it are introduced in Section~\ref{sec:bg:moo}.

\subsection{Formal Problem Definition}
\label{subsec:bg:prpa:formal}
A performance measure is formalized as a \emph{Key Performance Indicator}
(KPI) function $F : \mathcal{E}^{*} \times \mathbb{N} \to \mathbb{R}$ that,
given a trace $\sigma$ and a position $1 \le i \le |\sigma|$, returns the
value $F(\sigma, i)$ of that indicator for $\sigma$ after its first $i$
events~\cite{le2025balancing}. Evaluating a KPI at an intermediate position,
and not only on the completed trace, is what makes it usable while a case is
still running.

The prescription itself is a \emph{next step}: a pair
$(a, r) \in \mathcal{A} \times \mathcal{R}$ consisting of the activity to
perform next and the resource assigned to it. Given a running execution
$\sigma' = \langle e_1, \dots, e_k \rangle$, the prescriptive problem is to
choose the next event $e_{k+1}$ --- that is, the pair $(a, r)$ --- so that the
resulting completed trace
$\sigma_{\mathrm{pres}} = \langle e_1, \dots, e_k, e_{k+1}, \dots, e_n
\rangle$ attains the best possible KPI value. The candidate pairs are not
arbitrary: they are restricted to the activities that can actually follow the
current state of the case and to the resources able to perform them, a
feasibility constraint that is detailed, together with the rest of the
framework, in the methodology chapter. When a single KPI $F$ is considered,
``best'' simply means maximizing (or minimizing) $F$; when several KPIs
$F_1, \dots, F_n$ matter at once, they are aggregated into a single
\emph{utility function}
\begin{equation}
U_{\mathrm{KPI}}(\sigma, \boldsymbol{\lambda}, i)
  = \sum_{j=1}^{n} \lambda_j \, F_j(\sigma, i),
\label{eq:bg:prpa:utility}
\end{equation}
where $\boldsymbol{\lambda} = (\lambda_1, \dots, \lambda_n) \in [0, 1]^n$,
with $\sum_{j} \lambda_j = 1$, weights the relative importance of each
indicator according to the decision-maker's
priorities~\cite{le2025balancing}. The limitations of this linear
aggregation, and the alternative of not aggregating at all, are discussed in
Section~\ref{sec:bg:moo}.

\subsection{Key Performance Indicators}
\label{subsec:bg:prpa:kpi}
The objective of PrPA is stated informally as reaching a ``desired
outcome'', and a KPI is what turns that vague notion into a quantity that can
be optimized. Which quantity is appropriate depends entirely on the process
and on the organization's goals; recurring dimensions are the throughput time
of a case, its cost, some notion of quality or of regulatory compliance, and
the efficiency with which resources are used.

Two examples make the mapping concrete. In a loan-application process, a bank
that wants to grant more loans can take the \emph{acceptance rate} --- the
fraction of cases that reach a positive decision --- as a KPI to maximize,
whereas a bank that wants faster decisions can take the \emph{time to
decision} as a KPI to minimize. In a hospital, management may want to shorten
the \emph{length of stay} of patients in order to free capacity, while also
keeping the \emph{30-day readmission rate} low, so that discharging patients
earlier does not degrade the quality of care.

The two hospital KPIs illustrate a recurring difficulty: performance measures
often \emph{conflict}, so that improving one worsens another, and the right
balance between them is not fixed --- it depends on the domain, on the
regulatory context and on situational priorities. A prescription that is
optimal for one KPI can therefore be a poor choice overall. Making these
trade-offs explicit, instead of hiding them inside a single composite score,
is the problem addressed in Section~\ref{sec:bg:moo}.

% -----------------------------------------------------------------------------
% Planning notes -- material relocated to the methodology chapter, kept here
% for reference.
% -----------------------------------------------------------------------------
% Feasibility module:
% - Purpose: restrict the space to what is practically executable.
% - Activity filter: A|sigma(k) = activities that in the log follow the
%   abstract state of the prefix at least once (transition system, window w).
%   Fall back to progressively shorter suffixes if the exact prefix was never
%   seen.
% - Resource filter: R_a = resources that executed a at least once in the log.
% - Feasible prescription space: P_sigma(k) = {(a, r) : a in A|sigma(k),
%   r in R_a}.
% - Exclusion of "forbidden" activities (those defining the outcome).
%
% Search-space exploration strategies:
% - Exhaustive search: evaluates all feasible pairs; optimal w.r.t. the model
%   but costly when |P_sigma(k)| is large (many resources).
% - Heuristic / metaheuristic (genetic) methods: explore a subset.
%   (Algorithm detail belongs to Section \ref{sec:bg:moo} / methodology.)
%
% Extended single-KPI survey (if a separate Related Work chapter is wanted):
% - Resource allocation for global KPI improvement \cite{meneghello2024}.
% - de Leoni et al. on guidance for process continuation \cite{deleoni2020}.


% #############################################################################
\section{Conflicting KPIs and Multi-Objective Optimization}
\label{sec:bg:moo}
% #############################################################################
% Goal: the tools to understand why a Pareto front is used, how it is explored
% (enumeration, NSGA-II) and how a diversified subset of solutions is chosen.
% References: \cite{marler2004} (multi-objective survey); \cite{deb2002}
% (NSGA-II); \cite{mitchell1998} (genetic algorithms); \cite{mothilal2020}
% (DiCE); \cite{buliga2025}; p-dispersion / facility-location literature.
% Do NOT put here: the layout of the recommendations_*.csv files and the
% _objectives.csv sidecars -> Methodology / Evaluation.
% Bridge: "The quality of these decisions depends entirely on the quality -
% and the reliability - of the model that estimates the three objectives."

So far, most prescriptive-monitoring approaches optimize a single KPI. In
practice this is rarely enough: a business process is judged on several
performance measures at once, these measures are usually interdependent, and
improving one of them often comes at the expense of another. What is needed is
a way to reason about the whole set of measures together and to choose a
deliberate trade-off among them.

\subsection{Conflicting KPIs}
\label{subsec:bg:moo:devil}
The tension is easiest to see on an example. A hospital may want to shorten
the length of stay of its patients, in order to free capacity, and at the
same time keep the 30-day readmission rate low; discharging patients sooner
helps the first measure but tends to hurt the second, so no policy is best on
both counts and management has to settle on an acceptable compromise. The
same situation recurs across domains, and it is not limited to two measures:
the classical \emph{Devil's Quadrangle} already identifies four dimensions of
process performance --- time, cost, quality and flexibility --- that
typically cannot all be improved at once~\cite{dumas2018}.

\subsection{Scalarization and Its Limitations}
\label{subsec:bg:moo:scalarization}
One way to keep the convenience of a single optimum is to collapse the
several KPIs into one number. This is what the utility function of
Equation~\eqref{eq:bg:prpa:utility} does: a weighted sum of the individual
indicators, with weights $\boldsymbol{\lambda}$ chosen by the
decision-maker~\cite{marler2004}. Scalarization of this kind is simple and it
always yields one well-defined solution, but it has three drawbacks. It
forces the weights to be fixed \emph{before} any solution is seen, even
though the trade-off they encode is exactly what the decision-maker would
like to inspect; it returns a single point for each weight vector, hiding the
shape of the trade-off; and it requires the KPIs to be expressed on
comparable scales, which is why the time target is normalized as described in
Section~\ref{subsec:bg:ppm:normalization}.

\subsection{Pareto Optimality}
\label{subsec:bg:moo:pareto}
When the objectives genuinely conflict there is no single solution that is
best on all of them, and the problem instead has a whole set of equally
defensible compromises. This is captured by the notion of \emph{Pareto
optimality}~\cite{pareto1906,miettinen1999}. Consider a minimization problem
with $m$ objective functions $f_1, \dots, f_m$ and a decision vector $x$
ranging over a feasible set $X$. A solution $x_1$ is said to \emph{dominate}
a solution $x_2$, written $x_1 \prec x_2$, when it is no worse on every
objective and strictly better on at least one:
\begin{equation}
x_1 \prec x_2
\iff
\bigl( f_i(x_1) \le f_i(x_2)\ \ \forall\, i \in \{1, \dots, m\} \bigr)
\ \wedge\
\bigl( \exists\, j : f_j(x_1) < f_j(x_2) \bigr).
\label{eq:bg:moo:dominance}
\end{equation}
A solution is \emph{Pareto-optimal} --- equivalently \emph{efficient} or
\emph{non-dominated} --- if no feasible solution dominates it. The image of
the Pareto-optimal set in objective space is the \emph{Pareto front}
\begin{equation}
P^{*} = \bigl\{\, f(x) \;\big|\; x \in X,\ \nexists\, x' \in X :
f(x') \prec f(x) \,\bigr\}.
\label{eq:bg:moo:front}
\end{equation}
A point lies on the front simply because nothing dominates it; it need not
itself dominate any other point. The best and the worst value that each
objective attains over $P^{*}$ define the \emph{ideal} and the \emph{nadir}
point, which are convenient references for normalizing the objectives and for
measuring how close a given solution is to the (generally unattainable)
ideal. When two entire fronts have to be compared rather than two individual
points, a common scalar summary is the \emph{hypervolume}: the volume of
objective space that a front dominates with respect to a fixed reference
point~\cite{deb2001}.

Compared with scalarization, computing the Pareto front defers the trade-off
decision. The optimizer returns all the non-dominated compromises, and the
selection of one of them is made afterwards, with the cost of each
alternative in full view~\cite{honegumi}. In this thesis the front is built
on three objectives: the probability of a positive outcome, to be maximized;
the predicted remaining time, to be minimized; and the predictive standard
deviation of that time estimate, also to be minimized. The third objective
treats the model's own \emph{confidence} as a KPI, a choice motivated in
Section~\ref{sec:bg:uncertainty}.

\subsection{Front-Exploration Algorithms}
\label{subsec:bg:moo:algorithms}
% - Exact enumeration + front extraction: exact on small discrete spaces.
% - GENETIC ALGORITHMS: population, fitness, selection, crossover, mutation,
%   elitism, stopping criterion.
%    * The original paper's "custom fitness" approach: validity (val) +
%      alpha * proximity (prox), inspired by \cite{buliga2025}; DiCE as the
%      library \cite{mothilal2020}.
%    * NSGA-II \cite{deb2002}: non-dominated sorting into fronts, crowding
%      distance for diversity, tournament selection, elitism.
%      Implemented with pymoo (an integer variable indexing the (a,r) pair,
%      SBX + polynomial mutation + rounding repair), currently "parked"
%      because it is slower on small candidate sets: explain it anyway, it is
%      a methodological contribution.
% - Enumeration vs metaheuristic comparison: quality vs cost; robustness to
%   predictor errors on rare pairs (the paper's hypothesis: proximity favors
%   frequent pairs, which are estimated better).

\subsection{Selecting a Representative Subset of the Front}
\label{subsec:bg:moo:dispersion}
For a discrete problem the Pareto front can contain a large number of points,
some of them very close to one another. A recommender that has to present
alternatives to a user is often better served by a handful of clearly
distinct options than by the entire front. Choosing the $k$ most
representative points can be framed as a \emph{max-min dispersion}, or
\emph{$p$-dispersion}, problem: given $n$ candidate points with pairwise
distances $d_{ij}$, select a subset $S$ of size $k$ that maximizes the
smallest distance between any two selected points,
\begin{equation}
\max_{\substack{S \subseteq \{1, \dots, n\} \\ |S| = k}}\ \
\min_{\substack{i, j \in S \\ i \neq j}}\ d_{ij}.
\label{eq:bg:moo:dispersion}
\end{equation}
The problem was first studied in facility location, where the aim is to place
$p$ facilities so that the closest pair is as far apart as
possible~\cite{kuby1987}, and it has since been used more generally to
extract a maximally spread-out subset from a set of
candidates~\cite{maliszewski2012}. Because it is stated purely in terms of a
matrix of pairwise distances, the distance metric is a modeling choice; here
the natural one is the Euclidean distance between points in the normalized
objective space. The problem admits an exact mixed-integer programming
formulation, for which off-the-shelf solvers are
available~\cite{feng2022spopt}; the formulation actually used, together with
the rule that picks a single final recommendation, is given in the
methodology chapter.

Every objective on the front is produced by a predictive model, so the value
of the whole procedure rests on how accurate --- and how well calibrated ---
that model is. The next two sections describe it.


% #############################################################################
\section{Gradient Boosting and CatBoost}
\label{sec:bg:gbm}
% #############################################################################
% Goal: the model class used for the two predictors and what hyperparameters
% are / what role they play, plus their optimization. NO concrete values here.
% References: \cite{friedman2001} (gradient boosting); \cite{prokhorenkova2018}
% (CatBoost); \cite{akiba2019} (Optuna); \cite{bergstra2011} (TPE);
% \cite{galanti2023}.
% Do NOT put here: the search_spaces of 2_training_predictive_model.py, the
% number of trials, the best_hyperparams_*.json -> Experimental Setup (with a
% pointer to \ref{sec:bg:gbm}).
% Bridge: "A point model does not say how sure it is of what it predicts:
% uncertainty must be quantified."

\subsection{Tree Ensembles and Boosting}
\label{subsec:bg:gbm:boosting}
% - Decision / regression tree: splits, leaves, depth.
% - Bagging vs boosting (general idea).
% - Gradient boosting: sequential addition of trees fitting the negative
%   gradient of the loss; learning rate as shrinkage; number of trees
%   (iterations) and overfitting.

\subsection{CatBoost}
\label{subsec:bg:gbm:catboost}
% - Motivation for the choice: high predictive quality with contained
%   training time; native handling of categorical features; used in prior PPM
%   work \cite{galanti2023}.
% - Distinctive elements:
%    * ordered boosting / ordered target statistics (mitigates target leakage
%      / prediction shift);
%    * oblivious (symmetric) trees: same split at every level -> balanced
%      trees, fast inference, regularizing effect;
%    * growth alternatives (Depthwise, Lossguide) and min_data_in_leaf.
% - Losses used in the thesis (as CONCEPTS):
%    * Logloss for outcome classification;
%    * RMSE for regression; and the RMSEWithUncertainty variant, which
%      predicts mean AND variance (bridge to \ref{sec:bg:uncertainty}).
% - scale_pos_weight / class weights for imbalance (the positive-case
%   percentages differ a lot across logs: e.g. BAC ~82.5%, BPIC2012 ~47.9%).

\subsection{Hyperparameters: What They Are and What They Control}
\label{subsec:bg:gbm:hyperparams}
% - Definition: a parameter NOT learned from the data, fixed before training,
%   governing model capacity and the optimization process.
% - Grouped by ROLE (explain the role, not the value):
%    * Capacity / complexity: depth, grow_policy, min_data_in_leaf.
%    * Learning step: learning_rate, iterations.
%    * Regularization: l2_leaf_reg.
%    * Stochasticity / sub-sampling: bootstrap_type (Bayesian/Bernoulli/MVS),
%      bagging_temperature, subsample, colsample_bylevel.
%    * Imbalance: scale_pos_weight.
% - The bias-variance trade-off is governed by these groups.

\subsection{Hyperparameter Optimization}
\label{subsec:bg:gbm:hpo}
% - Why not by hand: large space, non-linear interactions.
% - Validation: hold-out, k-fold cross-validation, and - for processes - a
%   TEMPORAL split (avoid leakage from the future). Note the internal
%   "running" validation (the most recent training cases as validation).
% - Early stopping: stop adding trees when the validation metric stops
%   improving.
% - Search: grid / random search vs BAYESIAN OPTIMIZATION.
%    * Optuna: TPE (Tree-structured Parzen Estimator) as the surrogate;
%      define-by-run search space.
%    * Pruning (MedianPruner): interrupt unpromising trials.
% - Selection of the winning trial on ONE consistent metric (Logloss / RMSE)
%   and a final refit on 100% of the training data with the chosen tree count.


% #############################################################################
\section{Predictive Uncertainty Quantification}
\label{sec:bg:uncertainty}
% #############################################################################
% Goal: this is the THEORETICAL CORE of the contribution. Taxonomy of
% uncertainty, its estimation in ensembles and in gradient boosting,
% calibration. It must be the most developed section of the chapter.
% References: \cite{malinin2021} (Uncertainty in Gradient Boosting via
% Ensembles, ICLR 2021); \cite{hullermeier2021} (aleatoric vs epistemic,
% survey); \cite{lakshminarayanan2017} (deep ensembles); \cite{guo2017}
% (temperature scaling); \cite{kuleshov2018} (calibrated regression);
% \cite{angelopoulos2023} (conformal prediction); \cite{gal2016} (MC-dropout).
% Do NOT put here: the UncertaintyRegressor / UncertaintyClassifier classes,
% how the wrappers plug into the Pipeline, how out["std"] becomes objective #3
% -> Methodology. The obtained ECE/coverage/T/s values -> Evaluation.
% Bridge: "The Pareto front and the recommendations are evaluated with a
% process simulation, because a real A/B test is not possible."

\subsection{Taxonomy of Uncertainty}
\label{subsec:bg:uncertainty:taxonomy}
% - ALEATORIC uncertainty (data uncertainty): intrinsic / irreducible noise in
%   the data-generating process; does not shrink with more data.
%    * homoscedastic vs heteroscedastic.
% - EPISTEMIC uncertainty (knowledge / model uncertainty): the model's
%   ignorance due to lack of data / coverage; reducible with more data; high
%   on rare or out-of-distribution inputs.
% - TOTAL predictive uncertainty: the combination of the two.
% - Why distinguish them in the thesis: epistemic flags out-of-domain inputs /
%   rarely observed (a,r) pairs; aleatoric says how noisy the estimate is even
%   in the ideal case.
% - RESULT PREVIEW: on one-hot tabular data the regressor's epistemic part is
%   very small (~0.1% of the variance) and does NOT rise on never-seen inputs
%   (an unknown category -> an all-zero row every sub-ensemble agrees on). It
%   is reported anyway: the full decomposition is a result in itself
%   ("epistemic is negligible here, and why").

\subsection{Ensemble Methods for Uncertainty}
\label{subsec:bg:uncertainty:ensembles}
% - Idea: train/extract several models; the DISAGREEMENT among members
%   estimates the epistemic part, the AVERAGE member variance estimates the
%   aleatoric part.
% - Deep ensembles, MC-dropout, bootstrap ensembles: cite as context.
% - LAW OF TOTAL VARIANCE (regression):
%     Var[y] = E_m[ Var(y | m) ]  +  Var_m[ E(y | m) ]
%              \____ aleatoric ___/    \___ epistemic ___/
% - ENTROPY decomposition (classification):
%     H[ E_m p_m ]            = total uncertainty
%     E_m[ H(p_m) ]           = aleatoric uncertainty (expected entropy)
%     MI = total - aleatoric  = epistemic uncertainty (mutual information, BALD)
%   All in nats.

\subsection{Uncertainty in Gradient Boosting}
\label{subsec:bg:uncertainty:gbm}
% (\cite{malinin2021}.)
% - SGLB - Stochastic Gradient Langevin Boosting: noise injected into the
%   updates -> the model (approximately) samples from a posterior over the
%   parameters. In CatBoost: posterior_sampling=True.
% - VIRTUAL ENSEMBLES: from a single SGLB model several "sub-models" are
%   obtained using growing prefixes of the tree sequence -> an ensemble at the
%   cost of a single training run. Parameter: virtual_ensembles_count.
% - Regression with the RMSEWithUncertainty loss: each sub-model predicts
%   (mu_m, sigma_m^2); via the law of total variance:
%     data_std      = sqrt( mean_m sigma_m^2 )     (aleatoric)
%     knowledge_std = sqrt( var_m mu_m )           (epistemic)
%     total_std     = sqrt( data_var + knowledge_var )
% - Classification: prediction_type="TotalUncertainty" returns the expected
%   entropy (aleatoric) and the total entropy; the mutual information
%   (epistemic) is the difference.

\subsection{Calibration}
\label{subsec:bg:uncertainty:calibration}
% - What a CALIBRATED (probabilistic) model is: among the events predicted
%   with probability p, roughly a fraction p actually occur.
% - Diagnostics:
%    * reliability diagram;
%    * ECE - Expected Calibration Error: |confidence - accuracy| averaged over
%      probability bins;
%    * for regression: interval COVERAGE (fraction of |y - mu| within
%      +-1 sigma / +-2 sigma; target ~68% / ~95% for a Gaussian) and
%      SHARPNESS (typical / median interval width).
% - POST-HOC correction methods:
%    * TEMPERATURE SCALING \cite{guo2017}: a single scalar T,
%      p_cal = sigmoid( logit(p) / T ); T > 1 "softens" over-confident
%      probabilities. Monotone -> it does not reorder candidates by
%      probability, but it changes the magnitudes feeding the normalization
%      and the front distances.
%    * SIGMA RECALIBRATION for regression: a scalar factor s that rescales
%      every standard-deviation component so the empirical coverage matches
%      the nominal one; estimated on a held-out slice with a robust
%      median-of-standardized-residual estimator (avoids blow-ups when
%      CatBoost predicts a near-zero variance).
% - CONFORMAL PREDICTION as the RIGOROUS alternative \cite{angelopoulos2023}:
%   split-conformal on a dedicated holdout -> intervals with a distribution-
%   free marginal coverage guarantee. Cite it as "the rigorous alternative"
%   and motivate why the thesis used post-hoc recalibration instead (the
%   models are trained once, cost/benefit trade-off).

\subsection{Supporting Statistical Tools}
\label{subsec:bg:uncertainty:stats}
% (Compact recap.)
% - Variance and standard deviation as dispersion measures.
% - Shannon entropy of a Bernoulli; entropy as an uncertainty measure.
% - Mutual information / information gain.
% - The normal distribution, the 68-95-99.7 rule, the standardized residual,
%   the 0.6745 quantile of the normal (median absolute value).


% #############################################################################
\section{Business Process Simulation}
\label{sec:bg:simulation}
% #############################################################################
% Goal: what process simulation is, how a simulation model is obtained from
% the data, and with which metrics its fidelity is assessed.
% References: \cite{vinci2025prosit} (ProSiT); \cite{demoor2025} (SimBank);
% \cite{chapelacampa2024} (extraneous delays); \cite{vanderaalst2016} (Petri
% nets / soundness); \cite{vinci2025online} (online discovery).
% Do NOT put here: the prosit_simulation_results/... paths, the
% build_recommender_df code, the specific case ids -> Methodology / Evaluation.

\subsection{Why Simulation}
\label{subsec:bg:simulation:why}
% - Impossibility of an A/B test: one cannot intervene on a real process just
%   to evaluate.
% - COUNTERFACTUAL evaluation: simulate "what would have happened" following /
%   not following the recommendation.
% - Standard in PrPA for evaluating recommendations \cite{demoor2025}.

\subsection{Petri Nets and Process Models}
\label{subsec:bg:simulation:petri}
% - Petri net: places, transitions, arcs, marking, firing of a transition.
% - Workflow net and soundness (every case can always complete).
% - Initial / final marking; relation to the start/end of a case.
% - Replay and ALIGNMENT of a log prefix against the model: reconstruct the
%   marking reached by a historical prefix (log move, model move,
%   model-inserted activities).

\subsection{Stochastic Discrete-Event Simulation}
\label{subsec:bg:simulation:des}
% - Components of a process simulation model:
%    * case ARRIVAL process;
%    * transition branching weights / probabilities;
%    * activity EXECUTION times;
%    * WAITING times / queues;
%    * resource CALENDARS and availability;
%    * (extraneous delays \cite{chapelacampa2024}, as context).
% - Repeating the runs to mitigate stochasticity (e.g. 10 runs per scenario).

\subsection{Discovery of the Simulation Model}
\label{subsec:bg:simulation:discovery}
% - The simulation parameters are LEARNED from the historical log (resource
%   discovery, transition weights, time distributions, calendars).
% - METHODOLOGICAL NOTE: the simulator parameters are discovered from the FULL
%   log, not from the training split.
% - The Petri-net structure is a separate input (control-flow discovery),
%   distinct from the parameter discovery.
% - Online discovery for evolving processes \cite{vinci2025online} - context.

\subsection{ProSiT}
\label{subsec:bg:simulation:prosit}
% - The simulator used \cite{vinci2025prosit}: interactive and transparent
%   simulations.
% - What it returns: simulated logs; diagnostics (unreachable recommendations,
%   truncated "runaway" cases, non-fitting prefixes, model-inserted
%   activities).

\subsection{Assessing the Simulator's Fidelity}
\label{subsec:bg:simulation:fidelity}
% (STATISTICS.)
% - Real vs simulated distribution comparison:
%    * case-duration distributions (histograms / KDE, median);
%    * rate of positive-outcome cases, real vs simulated.
% - Distribution of the simulator's residuals.
% - (Optional) distribution-distance metrics: EMD / Wasserstein, KS; or a
%   comparison of summary statistics.

\subsection{KPI Improvement Metrics}
\label{subsec:bg:simulation:kpimetrics}
% (Definition; the numbers are in Evaluation.)
% - Delta_CO = average over prefixes of ( F_CO(sigma_pres) - F_CO(sigma_GT) ).
% - Delta_RT = average relative improvement of the remaining time w.r.t. the
%   ground truth.
% - Baseline / Ground Truth: no recommendation applied.
% - (Optional) confidence intervals on the deltas via bootstrap; 95%
%   confidence bands in the plots as lambda varies.


% #############################################################################
\section{Notation}
\label{sec:bg:notation}
% #############################################################################
% OPTIONAL but much appreciated by advisors. A half-page table of symbols,
% consistent with those used in the rest of the thesis.
% Uncomment and complete the table below (requires \usepackage{booktabs}).
%
% \begin{table}[h]
%   \centering
%   \begin{tabular}{ll}
%     \toprule
%     Symbol & Meaning \\
%     \midrule
%     $\mathcal{L}$              & event log \\
%     $\sigma,\ \sigma_{(k)}$    & trace, $k$-prefix \\
%     $\mathcal{A},\ \mathcal{R}$ & sets of activities and resources \\
%     $act,\ res$                & projection functions \\
%     $e = (c, a, \mathbf{d})$   & event (case id, activity, payload) \\
%     $S^{w}(\sigma_{(k)})$      & state abstraction with window $w$ \\
%     $\mathcal{P}_{\sigma_{(k)}}$ & feasible prescription space \\
%     $F_{CO},\ F_{RT}$          & KPIs: case outcome, remaining time (norm.) \\
%     $U_{KPI}$                  & utility function / combined KPI \\
%     $\lambda$                  & KPI trade-off weight \\
%     $\phi$                     & predictive model (regressor) \\
%     $p,\ p_{cal}$              & raw / calibrated outcome probability \\
%     $T$                        & calibration temperature \cite{guo2017} \\
%     $s$                        & std recalibration factor \\
%     $\sigma_{data}$            & aleatoric uncertainty (std) \\
%     $\sigma_{know}$            & epistemic uncertainty (std) \\
%     $\sigma_{tot}$             & total predictive uncertainty \\
%     $\mathrm{MI}$              & mutual information (epistemic, classif.) \\
%     $k$                        & n. of diversified recommendations per case \\
%     $\tau,\ \alpha$            & validity threshold, proximity weight \\
%     \bottomrule
%   \end{tabular}
%   \caption{Notation used in the thesis.}
%   \label{tab:notation}
% \end{table}


% =============================================================================
%  STATISTICAL TOOLS MAP -> WHERE THEY GO  (note to self, not for the thesis)
% =============================================================================
%   z-score, sigmoid, min-max, invertibility             -> \ref{subsec:bg:ppm:normalization}
%   Bayesian optimization / TPE, CV, temporal split,
%     early stopping, pruning                            -> \ref{subsec:bg:gbm:hpo}
%   variance / std, Shannon entropy, mutual information,
%     law of total variance                              -> \ref{subsec:bg:uncertainty:ensembles} / \ref{subsec:bg:uncertainty:stats}
%   ECE, reliability diagram, temperature scaling,
%     coverage, sharpness, conformal prediction          -> \ref{subsec:bg:uncertainty:calibration}
%   68-95-99.7 rule, standardized residual,
%     0.6745 quantile                                    -> \ref{subsec:bg:uncertainty:stats}
%   distribution comparison (KDE, median, EMD/KS),
%     residuals                                          -> \ref{subsec:bg:simulation:fidelity}
%   bootstrap / 95% confidence bands                     -> \ref{subsec:bg:simulation:kpimetrics} and Evaluation
%
%   HYPERPARAMETERS appear in Background ONLY as concepts
%   (\ref{subsec:bg:gbm:hyperparams}). Chosen values, search space, number of
%   Optuna trials and best-params tables -> Experimental Setup, with the
%   sentence "the concepts are introduced in Section~\ref{sec:bg:gbm}".
%
% =============================================================================
%  FINAL CHECKLIST BEFORE HANDING IN THE CHAPTER
% =============================================================================
%  [ ] Every section ends with a "bridge" to the next one.
%  [ ] No concrete experimental value in Background (concepts only).
%  [ ] No description of YOUR classes/files/functions (-> Methodology).
%  [ ] Every concept introduced is later used in Methodology or Evaluation
%      (no "orphan" background).
%  [ ] \ref{sec:bg:uncertainty} is the most developed section: it is the
%      contribution.
%  [ ] The aleatoric/epistemic taxonomy is explained BEFORE SGLB / virtual
%      ensembles.
%  [ ] Notation table consistent with the symbols in the rest of the thesis.
%  [ ] Every subsection has at least one solid bibliographic reference.
%  [ ] Agreed with the advisors whether Related Work stays here
%      (\ref{subsec:bg:prpa:sota}) or is a separate chapter.
% =============================================================================
